% Options for packages loaded elsewhere
\PassOptionsToPackage{unicode}{hyperref}
\PassOptionsToPackage{hyphens}{url}
%
\documentclass[
]{article}
\usepackage{amsmath,amssymb}
\usepackage{iftex}
\ifPDFTeX
  \usepackage[T1]{fontenc}
  \usepackage[utf8]{inputenc}
  \usepackage{textcomp} % provide euro and other symbols
\else % if luatex or xetex
  \usepackage{unicode-math} % this also loads fontspec
  \defaultfontfeatures{Scale=MatchLowercase}
  \defaultfontfeatures[\rmfamily]{Ligatures=TeX,Scale=1}
\fi
\usepackage{lmodern}
\ifPDFTeX\else
  % xetex/luatex font selection
\fi
% Use upquote if available, for straight quotes in verbatim environments
\IfFileExists{upquote.sty}{\usepackage{upquote}}{}
\IfFileExists{microtype.sty}{% use microtype if available
  \usepackage[]{microtype}
  \UseMicrotypeSet[protrusion]{basicmath} % disable protrusion for tt fonts
}{}
\makeatletter
\@ifundefined{KOMAClassName}{% if non-KOMA class
  \IfFileExists{parskip.sty}{%
    \usepackage{parskip}
  }{% else
    \setlength{\parindent}{0pt}
    \setlength{\parskip}{6pt plus 2pt minus 1pt}}
}{% if KOMA class
  \KOMAoptions{parskip=half}}
\makeatother
\usepackage{xcolor}
\usepackage[margin=1in]{geometry}
\usepackage{longtable,booktabs,array}
\usepackage{calc} % for calculating minipage widths
% Correct order of tables after \paragraph or \subparagraph
\usepackage{etoolbox}
\makeatletter
\patchcmd\longtable{\par}{\if@noskipsec\mbox{}\fi\par}{}{}
\makeatother
% Allow footnotes in longtable head/foot
\IfFileExists{footnotehyper.sty}{\usepackage{footnotehyper}}{\usepackage{footnote}}
\makesavenoteenv{longtable}
\usepackage{graphicx}
\makeatletter
\def\maxwidth{\ifdim\Gin@nat@width>\linewidth\linewidth\else\Gin@nat@width\fi}
\def\maxheight{\ifdim\Gin@nat@height>\textheight\textheight\else\Gin@nat@height\fi}
\makeatother
% Scale images if necessary, so that they will not overflow the page
% margins by default, and it is still possible to overwrite the defaults
% using explicit options in \includegraphics[width, height, ...]{}
\setkeys{Gin}{width=\maxwidth,height=\maxheight,keepaspectratio}
% Set default figure placement to htbp
\makeatletter
\def\fps@figure{htbp}
\makeatother
\setlength{\emergencystretch}{3em} % prevent overfull lines
\providecommand{\tightlist}{%
  \setlength{\itemsep}{0pt}\setlength{\parskip}{0pt}}
\setcounter{secnumdepth}{-\maxdimen} % remove section numbering
\newlength{\cslhangindent}
\setlength{\cslhangindent}{1.5em}
\newlength{\csllabelwidth}
\setlength{\csllabelwidth}{3em}
\newlength{\cslentryspacingunit} % times entry-spacing
\setlength{\cslentryspacingunit}{\parskip}
\newenvironment{CSLReferences}[2] % #1 hanging-ident, #2 entry spacing
 {% don't indent paragraphs
  \setlength{\parindent}{0pt}
  % turn on hanging indent if param 1 is 1
  \ifodd #1
  \let\oldpar\par
  \def\par{\hangindent=\cslhangindent\oldpar}
  \fi
  % set entry spacing
  \setlength{\parskip}{#2\cslentryspacingunit}
 }%
 {}
\usepackage{calc}
\newcommand{\CSLBlock}[1]{#1\hfill\break}
\newcommand{\CSLLeftMargin}[1]{\parbox[t]{\csllabelwidth}{#1}}
\newcommand{\CSLRightInline}[1]{\parbox[t]{\linewidth - \csllabelwidth}{#1}\break}
\newcommand{\CSLIndent}[1]{\hspace{\cslhangindent}#1}
\ifLuaTeX
  \usepackage{selnolig}  % disable illegal ligatures
\fi
\IfFileExists{bookmark.sty}{\usepackage{bookmark}}{\usepackage{hyperref}}
\IfFileExists{xurl.sty}{\usepackage{xurl}}{} % add URL line breaks if available
\urlstyle{same}
\hypersetup{
  pdftitle={Chapter2},
  pdfauthor={AndreyPovaliy},
  hidelinks,
  pdfcreator={LaTeX via pandoc}}

\title{Chapter2}
\author{AndreyPovaliy}
\date{2023-09-29}

\begin{document}
\maketitle

\hypertarget{ux433ux43bux430ux432ux430-2.-ux43cux430ux442ux440ux438ux430ux43bux44b-ux438-ux43cux435ux442ux43eux434ux44b}{%
\section{ГЛАВА 2. МАТРИАЛЫ И
МЕТОДЫ}\label{ux433ux43bux430ux432ux430-2.-ux43cux430ux442ux440ux438ux430ux43bux44b-ux438-ux43cux435ux442ux43eux434ux44b}}

В настоящей главее буду подробно представлены данные о дизайне
исследования, подробное описание преддиагностического этапа, описания
используемых диагностических методов, регистрируемых данных и
классификаторов,

\hypertarget{ux43eux431ux449ux430ux44f-ux445ux430ux440ux430ux43aux442ux435ux440ux438ux441ux442ux438ux43aux430-ux438ux441ux441ux43bux435ux434ux43eux432ux430ux43dux438ux44f-ux438ux441ux441ux43bux435ux434ux43eux432ux430ux43dux438ux435-ux43cux43eux436ux43dux43e-ux445ux430ux440ux430ux43aux442ux435ux440ux438ux437ux43eux432ux430ux442ux44c-ux43aux430ux43a-ux43aux43bux438ux43dux438ux447ux435ux441ux43aux43eux435-ux438ux441ux441ux43bux435ux434ux43eux432ux430ux43dux438ux435-ux434ux438ux430ux433ux43dux43eux441ux442ux438ux447ux435ux441ux43aux438ux445-ux430ux441ux43fux435ux43aux442ux43eux432-ux440ux430ux43dux43dux43dux435ux439-ux434ux438ux430ux433ux43dux43eux441ux442ux438ux43aux438-ux440ux430ux43aux430-ux43cux43eux43bux43eux447ux43dux43eux439-ux436ux435ux43bux435ux437ux44b.}{%
\subsection{2.1 Общая характеристика исследования Исследование можно
характеризовать как клиническое исследование диагностических аспектов
раннней диагностики рака молочной
железы.}\label{ux43eux431ux449ux430ux44f-ux445ux430ux440ux430ux43aux442ux435ux440ux438ux441ux442ux438ux43aux430-ux438ux441ux441ux43bux435ux434ux43eux432ux430ux43dux438ux44f-ux438ux441ux441ux43bux435ux434ux43eux432ux430ux43dux438ux435-ux43cux43eux436ux43dux43e-ux445ux430ux440ux430ux43aux442ux435ux440ux438ux437ux43eux432ux430ux442ux44c-ux43aux430ux43a-ux43aux43bux438ux43dux438ux447ux435ux441ux43aux43eux435-ux438ux441ux441ux43bux435ux434ux43eux432ux430ux43dux438ux435-ux434ux438ux430ux433ux43dux43eux441ux442ux438ux447ux435ux441ux43aux438ux445-ux430ux441ux43fux435ux43aux442ux43eux432-ux440ux430ux43dux43dux43dux435ux439-ux434ux438ux430ux433ux43dux43eux441ux442ux438ux43aux438-ux440ux430ux43aux430-ux43cux43eux43bux43eux447ux43dux43eux439-ux436ux435ux43bux435ux437ux44b.}}

Обектом исследования являются новобразования молочной железы,
регистрируемые при дигностике с использованием ручного ультразвукого
исследования (УЗИ), маммографии (ММГ) и автоматизированного объемного
УЗИ сканирования молочных желез (ABUS). Предметом исследования является
изучение диагностической эффективности автоматизированного объемного УЗИ
сканирования молочных желез, а также анализ наиболее важных факторов как
преддиагностическом этапе, так и при выполнении конкретного
диагностического метода, в частности при выполнении автоматизированного
объемного УЗИ сканирования молочных желез. Дизайн исследования можно
охарактеризовать ретро-проспективное наблюдательное исследование.
Протокол настоящего исследования был одобрен на заседании локального
этического комитета СЗГМУ им. Мечникова №00 от ХХ.ХХ.ХХХХ года.

Всего в исследование вошло 2794 пациенток. Критерии включения были
пациенты женского пола в возрасте от 18 до 80 лет; женщина, посетившая
врача для обследования рака молочной железы, с отсутствием видимых
признаков рака молочной железы. Критериями исключения были: женщины,
которые были беременны, кормили грудью или планировали забеременеть;
оперативное лечение в анамнезе (лампэктомия, мастэктомия, увеличение
груди), Эксцизионная или чрескожная биопсия за последние 12
месяцев,пациентки получавшие лечение по поводу рака молочных желез за
последние 12 месяцев. В это исследование была включена одна медицинская
клиника- СМТ, состоящая из двух зданий -амбулаторно-поликлического
комплекса и хирургического корпуса со стационаром. Все женщины из
амбулаторного отделения были приглашены для участия в нашем
исследовании, участники подписали форму информированного согласия. Все
обследования были проведены медицинским персоналом с соответствующей
квалификацией. Всем провели клинический осмотр, пальпацию, собрали
информацию о социально-демографических данных и потенциальных факторах
риска РМЖ.

Все пациентки были разделены на 2 независимые друг от друга выборки по
принципу возраста: выборка пациенток с возрастом до 40 лет и выборка
пациенток с возрастом от 40 лет включительно. Принцип разделения по
возрасту был связан с различными принципами использования
диагностических методов в скрининге рака молочной железы. У пациенток
моложе 40 лет ММГ менее эффективна, так как низкая чувствительность
метода и основным методом диагности используется УЗИ метод (Mandelson et
al. 2000). У пациенток моложе от 40 лет метод ММГ является обязательным
в скрининге РМЖ (Narayan, Lee, and Lehman 2020). Более подробно об
используемых методах можно прочитать в разделе ``2.3 Описание
диагностических методов''.

В выборку до 40 лет вошло 1511 пацинеток. Эта выборка пациенток была
разделена на две группы A и B. В группу A вошло 724 пациенток и в это
группе скрининг проводился с использовонием мануальной УЗИ диагностики и
и автоматизированного объемного УЗИ сканирования молочных желез
(исследуемая группа). В группу B вошло 787 пациенток и в этой группе
скрининг проводился только с исползовонием ручного УЗИ диагностики
(котнтрольная группа). В выборку 40 лет и старше вошло 1283 пацинеток.
Эта выборка пациенток была разделена на две группы С и D. В группу С
вошло 655 пациенток и в это группе скрининг проводился с использовонием
ручного УЗИ, маммографии и автоматизированного объемного УЗИ
сканирования молочных желез (исследуемая группа). В группу D вошло 628
пациенток и в этой группе скрининг проводился только с исползовонием
ручного УЗИ и маммографии (котнтрольная группа). Подробное описание
выборок и групп представлено в разделе ``2.5 Описание исследуемых
групп''. Включение указанного количества пациенток в группы было
основано на расчете мощности, более подробно можно прочитать в разделе
``2.6 Статистический анализ данных''.

Диагностический этап, где проводилась непосредственно диагностическое
исследование с использованием методом ручного УЗИ исследования,
маммографии и автоматизированного объемного УЗИ сканирования молочных
желез в зависимости от группы распределения.

Ппроизводился сбор данных для последующего выления наиболее значимых
факторов, а также выявления эффективности изучаемых методов и
прогностической ценности метода.

\hypertarget{ux43eux43fux438ux441ux430ux43dux438ux435-ux434ux438ux430ux433ux43dux43eux441ux442ux438ux447ux435ux441ux43aux438ux445-ux43cux435ux442ux43eux434ux43eux432}{%
\subsection{2.2 Описание диагностических
методов}\label{ux43eux43fux438ux441ux430ux43dux438ux435-ux434ux438ux430ux433ux43dux43eux441ux442ux438ux447ux435ux441ux43aux438ux445-ux43cux435ux442ux43eux434ux43eux432}}

\hypertarget{ux440ux443ux447ux43dux43eux435-ux443ux437ux438-ux438ux441ux441ux43bux435ux434ux43eux432ux430ux43dux438ux435}{%
\subsubsection{2.2.1 Ручное УЗИ
исследование}\label{ux440ux443ux447ux43dux43eux435-ux443ux437ux438-ux438ux441ux441ux43bux435ux434ux43eux432ux430ux43dux438ux435}}

HHUS выполняли два врача ультразвуковой диагностики со стажем работы
более 7 лет. Исследование проводилось в положении лежа, с руками за
головой, с последовательным сканированием каждого квадранта обеих
молочных железе в сагитальной и аксиальной плоскостях, с исследованием
ретроареолярной области и аксиллярных областей с двух сторон.
Устройства, используемые для проведения HHUS включали GE LOGIQS 8 (GE
Medical Systems, Милуоки, Висконсин, США), Toshiba Aplio 300(Canon
Япония)- ультразвуковые системы экспертного класса.

\hypertarget{ux430ux432ux442ux43eux43cux430ux442ux438ux437ux438ux440ux43eux432ux430ux43dux43dux43eux435-ux43eux431ux44aux435ux43cux43dux43eux435-ux443ux437ux438-ux441ux43aux430ux43dux438ux440ux43eux432ux430ux43dux438ux435-ux43cux43eux43bux43eux447ux43dux44bux445-ux436ux435ux43bux435ux437-abus}{%
\subsubsection{2.3.2 Автоматизированное объемное УЗИ сканирование
молочных желез
(ABUS)}\label{ux430ux432ux442ux43eux43cux430ux442ux438ux437ux438ux440ux43eux432ux430ux43dux43dux43eux435-ux43eux431ux44aux435ux43cux43dux43eux435-ux443ux437ux438-ux441ux43aux430ux43dux438ux440ux43eux432ux430ux43dux438ux435-ux43cux43eux43bux43eux447ux43dux44bux445-ux436ux435ux43bux435ux437-abus}}

D-автоматизированная ультразвуковая система Invenia (ABUS) оцениваемая в
данном исследовании, производства GE Healthcare (Саннивейл, Калифорния,
США) 2018 года выпуска--- это компьютерная система для оценки плотной
молочной железы.

Каждая молочная железа была визуализирована в трех проекциях: боковой
(LAT), переднезадней (AP).и медиальный (MED) с автоматическим датчиком с
линейной матрицей от 6 до 14 МГц, прикрепленным к жесткой компрессионной
пластине.(площадь 15,4×17,0×5,0 см).Индивидуальная мембрана для датчика
была использована для каждой пациентки. Во время исследования система
получала до 300 2D-срезов и реконструировала их в коронарной плоскости.
Стандартизированный процесс осмотра включает использование
запатентованной коронарной плоскости для быстрой навигации по молочной
железе, а также использование «режима обзора», позволяющего врачу быстро
интерпретировать изображения. Время сбора данных для каждой проекции
составляло 60 с, примерно по 3--4 мин на каждую молочную железу.
Обследование проводили в положении лежа. Полотенце было помещено под
плечом, что помогло расправить ткань молочной железы равномерно, соском
к потолку. На молочные железы равномерно наносили гипоаллергенный
ультразвуковой гель с дополнительным количество на область соска..
Возможно применение трех уровней компрессии датчика для исследования
молочных желез для получения наилучшей визуализации с учетом комфорта
пациента.

Сканирование ABUS было непрерывным и автоматизированным.В течение
исследования женщин попросили не двигаться, не разговаривать и дышать
ровно. Выполнял исследование сертифицированный персонал со средним
медицинским образованием. После завершения сбора данных ультразвуковой
системой весь массив передавался на специальную рабочую станцию для
интерпретации. Оценку изображений ABUS выполнял один врач ультразвуковой
диагностики, со стажем работы более 7 лет. Общее время, необходимое для
подготовки пациента и получения ABUS, фиксировалось, в каждом случае и
варьировалось примерно между 10 и 15 мин.

\hypertarget{ux43cux430ux43cux43cux43eux433ux440ux430ux444ux438ux44f}{%
\subsubsection{2.3.3
Маммография}\label{ux43cux430ux43cux43cux43eux433ux440ux430ux444ux438ux44f}}

Пациентки старше 40 лет прошли двухпроекционную цифровую маммографию (в
медиалатеральной косой и краниокаудальной проекции)обеих молочных желез.
Маммографию также выполняли женщинам моложе 40 лет в случае
положительного семейного или личного анамнеза - рак молочной железы.
Использумемое оборудование- Planmed Clarity 3D с функцией томосинтеза
(Финляндия). Оценку изображений проводил один рентгенолог со стажем
работы более 10 лет. \#\#\# 2.3.4 Магнитно-резонанская томография

\hypertarget{ux43eux446ux435ux43dux43aux430-ux432ux44bux44fux432ux43bux435ux43dux43dux44bux445-ux438ux437ux43cux435ux43dux435ux43dux438ux439.}{%
\subsubsection{2.3.5 Оценка выявленных
изменений.}\label{ux43eux446ux435ux43dux43aux430-ux432ux44bux44fux432ux43bux435ux43dux43dux44bux445-ux438ux437ux43cux435ux43dux435ux43dux438ux439.}}

Результаты выявленных изменений были классифицированы по BI-RADS
отсортированы по шести категориям (Spak et al. 2017): 0 = неполные
данные , 1 = нормальные, 2 = доброкачественные, 3 =вероятно
доброкачественные, 4 = подозрительные, 5 = с высокой степенью
вероятности злокачественные, для HHUS, ABUS и MG. Для категории Birads 3
обязательно проводилась магнитно-резонансная томография (МРТ) для
выявления истинно отрицательных и ложноотрицательных результатов.
Женщинам, отнесенным к категории 4 или 5, проводилась трепан биопсия под
уз-наведением или стереотаксическим наведением с последующей
морфологической и при необходимости иммуногистохимической верификацией.

\hypertarget{ux43eux43fux438ux441ux430ux43dux438ux435-ux440ux435ux433ux438ux441ux442ux440ux438ux440ux443ux435ux43cux44bux445-ux434ux430ux43dux43dux44bux445-ux438-ux438ux445-ux43aux43bux430ux441ux441ux438ux444ux438ux43aux430ux442ux43eux440ux432}{%
\subsection{2.4 Описание регистрируемых данных и их
классификаторв}\label{ux43eux43fux438ux441ux430ux43dux438ux435-ux440ux435ux433ux438ux441ux442ux440ux438ux440ux443ux435ux43cux44bux445-ux434ux430ux43dux43dux44bux445-ux438-ux438ux445-ux43aux43bux430ux441ux441ux438ux444ux438ux43aux430ux442ux43eux440ux432}}

\hypertarget{ux43eux43fux438ux441ux430ux43dux438ux435-ux438ux441ux441ux43bux435ux434ux443ux435ux43cux44bux445-ux433ux440ux443ux43fux43f}{%
\subsection{2.5 Описание исследуемых
групп}\label{ux43eux43fux438ux441ux430ux43dux438ux435-ux438ux441ux441ux43bux435ux434ux443ux435ux43cux44bux445-ux433ux440ux443ux43fux43f}}

\hypertarget{ux43eux431ux449ux435ux435-ux43eux43fux438ux441ux430ux43dux438ux435-ux432ux44bux431ux43eux440ux43aux438-ux43fux430ux446ux438ux435ux43dux442ux43eux43a-ux434ux43e-40-ux43bux435ux442}{%
\subsubsection{2.5.1 Общее описание выборки пациенток до 40
лет}\label{ux43eux431ux449ux435ux435-ux43eux43fux438ux441ux430ux43dux438ux435-ux432ux44bux431ux43eux440ux43aux438-ux43fux430ux446ux438ux435ux43dux442ux43eux43a-ux434ux43e-40-ux43bux435ux442}}

С февраля 2019 по май 2023 года было исследовано 1511 пациенток с
возрастом до 40 лет.

Медиана возраста пациенток выборке до 40 лет составила 35 {[}Q1-Q3:
32;37{]} лет. Минимальный возраст составил 20 лет. Максимальный возраст
составил 39 лет. Медиана роста пациенток выборке до 40 лет составил 167
{[}Q1-Q3: 164;170{]} см. Медиана веса пациенток выборке до 40 лет
составил 59 {[}Q1-Q3: 55;65{]} кг.

В выборке до 40 лет были поставлены следующие диагнозы при первичном
осмотре: «Без патологии» было 29.91\% пациенток(452/1511 случаев)
{[}95\% ДИ 0.28;0.32{]}, «Рак» был поставлен в 0.79\% (12/1511 случаев)
{[}95\% ДИ 0;0.01{]}, «Фиброаденома» была поставлена в 19.99\% (302/1511
случаев) {[}95\% ДИ 0.18;0.22{]}, «Фиброзно-кистозная болезнь» в 25.61\%
(387/1511 случаев) {[}95\% ДИ 0.23;0.28{]}, «Киста» была в 6.42\%
(97/1511 случаев) {[}95\% ДИ 0.05;0.08{]}, «Листовидная опухоль» - 0.6\%
(9/1511 случаев) {[}95\% ДИ 0;0.01{]}, «Внутрипротоковая папиллома» -
0.4\% (6/1511 случаев) {[}95\% ДИ 0;0.01{]}, «Гиперплазия» было 0.4\%
(6/1511 случаев) {[}95\% ДИ 0;0.01{]}, «Диффузный фиброаденоматоз» -
15.42\% (233/1511 случаев) {[}95\% ДИ 0.14;0.17{]} и «Локализованный
фиброаденоматоз» - 0.46\% (7/1511 случаев) {[}95\% ДИ 0;0.01{]}.

Были выялены следующие жалобы: «выделения из соска» было 0.66\% (10/1511
случаев) {[}95\% ДИ 0;0.01{]}, боль была в 5.03\% (76/1511 случаев)
{[}95\% ДИ 0.04;0.06{]}, «уплотнение» было в 9.4\% (142/1511 случаев)
{[}95\% ДИ 0.08;0.11{]} и не было жалоб в 84.91\% (1283/1511 случаев)
{[}95\% ДИ 0.83;0.87{]}.

Репродуктивный возраст был в 96.96\% (1465/1511 случаев) {[}95\% ДИ
0.96;0.98{]}, пременопауза была в 3.04\% (46/1511 случаев) {[}95\% ДИ
0.02;0.04{]}.

Операций на молочной железе в анамнезе не было в 98.41\% (1487/1511
случаев) {[}95\% ДИ 0.98;0.99{]} и были в 1.59\% (24/1511 случаев)
{[}95\% ДИ 0.01;0.02{]} .

В выборке до 40 лет прием гормональных препаратов был в 21.05\%
(310/1511 случаев) {[}95\% ДИ 0.19;0.23{]} и не отмечался в 8.95\%
(1163/1511 случаев) {[}95\% ДИ 0.75;0.79{]}.

В выборке до 40 лет генетическая предрасположенность нет было в 68.57\%
(1023/1511 случаев) {[}95\% ДИ 0.65;0.7{]}, в то время как
предрасположенность отмечалась в 31.43\% (469/1511 случаев) {[}95\% ДИ
0.29;0.33{]}.

Была выявлена мутация BRCA BRCA1 было 0.4\% (6/1511 случаев) {[}95\% ДИ
0;0.01{]}, мутаций BRCA не выявлено было 1.19\% (18/1511 случаев)
{[}95\% ДИ 0.01;0.02{]}. В остальных случаях исследования мутаций не
проводилось.

Сторона поражения при осмотре была правая в 11.25\% (170/1511 случаев)
{[}95\% ДИ 0.1;0.13{]}, левая в 6.55\% (99/1511 случаев) {[}95\% ДИ
0.05;0.08{]}, поражения обеих сторон в 3.97\% (60/1511 случаев) {[}95\%
ДИ 0.03;0.05{]}, в остальных случаях образование не пальпировало -
78.23\% (1182/1511 случаев) {[}95\% ДИ 0.76;0.8{]}.

Кожные симптомы при осмотре нет было во всех 100\% (1511/1511 случаев)
{[}95\% ДИ 1;1{]}.

В выборке до 40 лет втягивание соска при осмотре не было в 98.81\%
(1493/1511 случаев) {[}95\% ДИ 0.98;0.99{]}, было отмечено в 1.19\%
(18/1511 случаев) {[}95\% ДИ 0.01;0.02{]}.

Выделения из соска при осмотре не было в 99.54\% (1504/1511 случаев)
{[}95\% ДИ 0.99;1{]}., кровянистые выделения были в 0.46\% (7/1511
случаев) {[}95\% ДИ 0;0.01{]}.

Тип плотности по ACR при осмотре был «С» в 70.11\% (1046/1511 случаев)
{[}95\% ДИ 0.67;0.72{]}, тип «D» был в 29.89\% (446/1511 случаев)
{[}95\% ДИ 0.27;0.32{]}.

Квадрант локализации верхне-нараружный был в 8.01\% (121/1511 случаев)
{[}95\% ДИ 0.07;0.1{]}, верхне-внутренний - 3.97\% (60/1511 случаев)
{[}95\% ДИ 0.03;0.05{]}, нижне-наружный - 2.38\% (36/1511 случаев)
{[}95\% ДИ 0.02;0.03{]}, нижне-внутренний было в 1.19\% (18/1511
случаев) {[}95\% ДИ 0.01;0.02{]}, граница внутреннего - 1.72\% (26/1511
случаев) {[}95\% ДИ 0.01;0.03{]}, граница наружного - 1.19\% (18/1511
случаев) {[}95\% ДИ 0.01;0.02{]}, граница верхнего - 2.91\% (44/1511
случаев) {[}95\% ДИ 0.02;0.04{]}, граница нижних - 0.86\% (13/1511
случаев) {[}95\% ДИ 0;0.02{]}, вся железа - 0.4\% (6/1511 случаев)
{[}95\% ДИ 0;0.01{]}.

\hypertarget{ux43eux446ux435ux43dux43aux430-ux43eux434ux43dux43eux440ux43eux434ux43dux43eux441ux442ux438-ux43fux430ux446ux438ux435ux43dux442ux43eux43a-ux434ux43e-40-ux43bux435ux442}{%
\subsubsection{2.5.2 Оценка однородности пациенток до 40
лет}\label{ux43eux446ux435ux43dux43aux430-ux43eux434ux43dux43eux440ux43eux434ux43dux43eux441ux442ux438-ux43fux430ux446ux438ux435ux43dux442ux43eux43a-ux434ux43e-40-ux43bux435ux442}}

В группу A вошло 787 пациенток. Медиана возраста составила 35 {[}Q1-Q3:
32;37{]} лет. Минимальный возраст составил 20 лет. Максимальный возраст
составил 39 лет. Медиана роста составила 167 {[}Q1-Q3: 164;170{]} см.
Медиана веса составила 58 {[}Q1-Q3: 55;65{]} кг. Пациенток в группе B в
исследование вошло 724 пациентов. Медиана возраста составила 34
{[}Q1-Q3: 31;37{]} лет. Минимальный возраст составил 20 лет.
Максимальный возраст составил 39 лет. Медиана роста оставила 167
{[}Q1-Q3: 164;170{]} см. Медиана веса составила 60 {[}Q1-Q3: 55;65{]}
кг. Статистически значимой разницы по возрасту, росту, весу не было
выявлено.

\emph{Дигноз при первичном осмотре}

В группе А были поставлены следующие диагнозы: «Без патологии» было
32.15\% (253/787 случаев ) {[}95\% ДИ 0.29;0.36{]}, «Рак» было 0.76\%
(6/787 случаев ) {[}95\% ДИ 0;0.02{]}, «Фиброаденома» в 15.63\% (123/787
случаев ) {[}95\% ДИ 0.13;0.18{]}, «Фиброзно-кистозная болезнь» в
28.72\% (226/787 случаев ) {[}95\% ДИ 0.26;0.32{]}, «Киста» в 4.45\%
(35/787 случаев ) {[}95\% ДИ 0.03;0.06{]},«Листовидная опухоль» в 0.38\%
(3/787 случаев ) {[}95\% ДИ 0;0.01{]}, «Диффузный фиброаденоматоз» в
17.03\% (134/787 случаев ) {[}95\% ДИ 0.15;0.2{]},«Локализованный
фиброаденоматоз» в 0.89\% (7/787 случаев ) {[}95\% ДИ 0;0.02{]}.

В группе B были поставлены диагнозы после осмотра: «Без патологии» было
27.49\% (199/724 случаев ) {[}95\% ДИ 0.24;0.31{]}, «Рак» было 0.83\%
(6/724 случаев ) {[}95\% ДИ 0;0.02{]},«Фиброаденома» в 24.72\% (179/724
случаев ) {[}95\% ДИ 0.22;0.28{]}, «Фиброзно-кистозная болезнь» в
22.24\% (161/724 случаев ) {[}95\% ДИ 0.19;0.25{]}, «Киста» в 8.56\%
(62/724 случаев ) {[}95\% ДИ 0.07;0.11{]}, «Листовидная опухоль» в
0.83\% (6/724 случаев ) {[}95\% ДИ 0;0.02{]}, «Внутрипротоковая
папиллома'' в 0.83\% (6/724 случаев ) {[}95\% ДИ 0;0.02{]},
«Гиперплазия» в 0.83\% (6/724 случаев ) {[}95\% ДИ 0;0.02{]}, «Диффузный
фиброаденоматоз» в 13.67\% (99/724 случаев ) {[}95\% ДИ 0.11;0.16{]}.

\includegraphics{images/Primary_diagnosis}

Разница между группами по первичному диагнозу была выявлена (p-value
\textless{} 0,01). Разница объясняется тем, что диагнозы
``Внутрипротоковая папиллома'', ``Гиперплазия'' в группе А не было
вылено, в то время как ``Локализованный фиброаденоматоз'' не был
зарегистрирован в группе В. Также в групе В вывлено больше пациенток с
диагнозом ``Фиброаденома'', а также ``Киста''. Учитывая, что наиболее
значимым диагнозом являлся ``Рак'', то пересчет с учетом приоритетной
значимости этого диагноза привел к тому, что группы распределеный
однородно (p-уровень = 0.77).

\begin{longtable}[]{@{}
  >{\raggedright\arraybackslash}p{(\columnwidth - 10\tabcolsep) * \real{0.1667}}
  >{\raggedright\arraybackslash}p{(\columnwidth - 10\tabcolsep) * \real{0.1667}}
  >{\raggedright\arraybackslash}p{(\columnwidth - 10\tabcolsep) * \real{0.1667}}
  >{\raggedright\arraybackslash}p{(\columnwidth - 10\tabcolsep) * \real{0.1667}}
  >{\raggedright\arraybackslash}p{(\columnwidth - 10\tabcolsep) * \real{0.1667}}
  >{\raggedright\arraybackslash}p{(\columnwidth - 10\tabcolsep) * \real{0.1667}}@{}}
\toprule\noalign{}
\begin{minipage}[b]{\linewidth}\raggedright
Диагноз
\end{minipage} & \begin{minipage}[b]{\linewidth}\raggedright
Процентная доля
\end{minipage} & \begin{minipage}[b]{\linewidth}\raggedright
95\% ДИ
\end{minipage} & \begin{minipage}[b]{\linewidth}\raggedright
Процентная доля
\end{minipage} & \begin{minipage}[b]{\linewidth}\raggedright
95\% ДИ
\end{minipage} & \begin{minipage}[b]{\linewidth}\raggedright
p-value
\end{minipage} \\
\midrule\noalign{}
\endhead
\bottomrule\noalign{}
\endlastfoot
& Группа А & & Группа В & & \textless{} 0,01 \\
Без патологии & 32.147 \% ( 253 / 787 случаев) & {[} 0.29 ; 0.36 {]} &
27.486 \% ( 199 / 724 случаев) & {[} 0.24 ; 0.31 {]} & \\
Рак & 0.762 \% ( 6 / 787 случаев) & {[} 0 ; 0.02 {]} & 0.829 \% ( 6 /
724 случаев) & {[} 0 ; 0.02 {]} & 0.77 \\
Саркома & 0 \% ( 0 / 787 случаев) & {[} 0 ; 0.01 {]} & 0 \% ( 0 / 724
случаев) & {[} 0 ; 0.01 {]} & \\
Фиброаденома & 15.629 \% ( 123 / 787 случаев) & {[} 0.13 ; 0.18 {]} &
24.724 \% ( 179 / 724 случаев) & {[} 0.22 ; 0.28 {]} & \\
Фиброзно-кистозная болезнь & 28.717 \% ( 226 / 787 случаев) & {[} 0.26 ;
0.32 {]} & 22.238 \% ( 161 / 724 случаев) & {[} 0.19 ; 0.25 {]} & \\
Киста & 4.447 \% ( 35 / 787 случаев) & {[} 0.03 ; 0.06 {]} & 8.564 \% (
62 / 724 случаев) & {[} 0.07 ; 0.11 {]} & \\
Листовидная опухоль & 0.381 \% ( 3 / 787 случаев) & {[} 0 ; 0.01 {]} &
0.829 \% ( 6 / 724 случаев) & {[} 0 ; 0.02 {]} & \\
Внутрипротоковая папиллома & 0 \% ( 0 / 787 случаев) & {[} 0 ; 0.01 {]}
& 0.829 \% ( 6 / 724 случаев) & {[} 0 ; 0.02 {]} & \\
Гиперплазия & 0 \% ( 0 / 787 случаев) & {[} 0 ; 0.01 {]} & 0.829 \% ( 6
/ 724 случаев) & {[} 0 ; 0.02 {]} & \\
Диффузный фиброаденоматоз & 17.027 \% ( 134 / 787 случаев) & {[} 0.15 ;
0.2 {]} & 13.674 \% ( 99 / 724 случаев) & {[} 0.11 ; 0.16 {]} & \\
Локализованный фиброаденоматоз & 0.889 \% ( 7 / 787 случаев) & {[} 0 ;
0.02 {]} & 0 \% ( 0 / 724 случаев) & {[} 0 ; 0.01 {]} & \\
Мастит & 0 \% ( 0 / 787 случаев) & {[} 0 ; 0.01 {]} & 0 \% ( 0 / 724
случаев) & {[} 0 ; 0.01 {]} & \\
\end{longtable}

\emph{Жалобы при первичном осмотре}

В группе А были выявлены следующие жалобы: выделения из соска в 0.38\%
(3/787 случаев ) {[}95\% ДИ 0;0.01{]}, боль была в 4.83\% (38/787
случаев ) {[}95\% ДИ 0.03;0.07{]}, уплотнение было в 8.01\% (63/787
случаев ) {[}95\% ДИ 0.06;0.1{]}, в 86.79\% (683/787 случаев ) {[}95\%
ДИ 0.84;0.89{]} не было жалоб.

В группе B были выявлены следующие жалобы: выделения из соска были в
0.97\% (7/724 случаев ) {[}95\% ДИ 0;0.02{]}, боль была в 5.25\% (38/724
случаев ) {[}95\% ДИ 0.04;0.07{]}, уплотнение было в 10.91\% (79/724
случаев ) {[}95\% ДИ 0.09;0.13{]}, жалоб не было в 82.87\% (600/724
случаев ) {[}95\% ДИ 0.8;0.86{]}.

\includegraphics{images/complaints}

Разница между группами по выявленным жалобам не было выявлено составила
уровень значимости составил (p-уровень = 0.21). По составу жалоб группы
распределены равномерно, можно отметить не значимую разницу в общем
количестве пациенток с отсутвием жалоб.

\begin{longtable}[]{@{}
  >{\raggedright\arraybackslash}p{(\columnwidth - 10\tabcolsep) * \real{0.1667}}
  >{\raggedright\arraybackslash}p{(\columnwidth - 10\tabcolsep) * \real{0.1667}}
  >{\raggedright\arraybackslash}p{(\columnwidth - 10\tabcolsep) * \real{0.1667}}
  >{\raggedright\arraybackslash}p{(\columnwidth - 10\tabcolsep) * \real{0.1667}}
  >{\raggedright\arraybackslash}p{(\columnwidth - 10\tabcolsep) * \real{0.1667}}
  >{\raggedright\arraybackslash}p{(\columnwidth - 10\tabcolsep) * \real{0.1667}}@{}}
\toprule\noalign{}
\begin{minipage}[b]{\linewidth}\raggedright
Жалобы
\end{minipage} & \begin{minipage}[b]{\linewidth}\raggedright
Процентная доля
\end{minipage} & \begin{minipage}[b]{\linewidth}\raggedright
95\% ДИ
\end{minipage} & \begin{minipage}[b]{\linewidth}\raggedright
Процентная доля
\end{minipage} & \begin{minipage}[b]{\linewidth}\raggedright
95\% ДИ
\end{minipage} & \begin{minipage}[b]{\linewidth}\raggedright
p-value
\end{minipage} \\
\midrule\noalign{}
\endhead
\bottomrule\noalign{}
\endlastfoot
& Группа А & & Группа Б & & 0.21 \\
нет & 86.785 \% ( 683 / 787 случаев) & {[} 0.84 ; 0.89 {]} & 82.873 \% (
600 / 724 случаев) & {[} 0.8 ; 0.86 {]} & \\
есть опухоль & 0 \% ( 0 / 787 случаев) & {[} 0 ; 0.01 {]} & 0 \% ( 0 /
724 случаев) & {[} 0 ; 0.01 {]} & \\
выделения из соска & 0.381 \% ( 3 / 787 случаев) & {[} 0 ; 0.01 {]} &
0.967 \% ( 7 / 724 случаев) & {[} 0 ; 0.02 {]} & \\
боль & 4.828 \% ( 38 / 787 случаев) & {[} 0.03 ; 0.07 {]} & 5.249 \% (
38 / 724 случаев) & {[} 0.04 ; 0.07 {]} & \\
уплотнение & 8.005 \% ( 63 / 787 случаев) & {[} 0.06 ; 0.1 {]} & 10.912
\% ( 79 / 724 случаев) & {[} 0.09 ; 0.13 {]} & \\
гиперемия & 0 \% ( 0 / 787 случаев) & {[} 0 ; 0.01 {]} & 0 \% ( 0 / 724
случаев) & {[} 0 ; 0.01 {]} & \\
гипертермия & 0 \% ( 0 / 787 случаев) & {[} 0 ; 0.01 {]} & 0 \% ( 0 /
724 случаев) & {[} 0 ; 0.01 {]} & \\
\end{longtable}

\emph{Репродуктивный статус}

В группе А репродуктивный возраст был в 96.7\% (761/787 случаев )
{[}95\% ДИ 0.95;0.98{]}, пременопауза было 3.3\% (26/787 случаев )
{[}95\% ДИ 0.02;0.05{]}.

В группе B репродуктивный статус был репродуктивный возраст был в
97.24\% (704/724 случаев ) {[}95\% ДИ 0.96;0.98{]}, пременопауза в
2.76\% (20/724 случаев ) {[}95\% ДИ 0.02;0.04{]}.

\includegraphics{images/satus_reproductive}

Разница между группами по репродуктивному статусу составила p-уровень=
0.55 . Группы распределены однородно.

\begin{longtable}[]{@{}
  >{\raggedright\arraybackslash}p{(\columnwidth - 10\tabcolsep) * \real{0.1667}}
  >{\raggedright\arraybackslash}p{(\columnwidth - 10\tabcolsep) * \real{0.1667}}
  >{\raggedright\arraybackslash}p{(\columnwidth - 10\tabcolsep) * \real{0.1667}}
  >{\raggedright\arraybackslash}p{(\columnwidth - 10\tabcolsep) * \real{0.1667}}
  >{\raggedright\arraybackslash}p{(\columnwidth - 10\tabcolsep) * \real{0.1667}}
  >{\raggedright\arraybackslash}p{(\columnwidth - 10\tabcolsep) * \real{0.1667}}@{}}
\toprule\noalign{}
\begin{minipage}[b]{\linewidth}\raggedright
Репродуктивный статус
\end{minipage} & \begin{minipage}[b]{\linewidth}\raggedright
Процентная доля
\end{minipage} & \begin{minipage}[b]{\linewidth}\raggedright
95\% ДИ
\end{minipage} & \begin{minipage}[b]{\linewidth}\raggedright
Процентная доля
\end{minipage} & \begin{minipage}[b]{\linewidth}\raggedright
95\% ДИ
\end{minipage} & \begin{minipage}[b]{\linewidth}\raggedright
p-value
\end{minipage} \\
\midrule\noalign{}
\endhead
\bottomrule\noalign{}
\endlastfoot
& Группа А & & Группа Б & & 0.55 \\
репродуктивный возраст & 96.696 \% ( 761 / 787 случаев) & {[} 0.95 ;
0.98 {]} & 97.238 \% ( 704 / 724 случаев) & {[} 0.96 ; 0.98 {]} & \\
менопауза до 5 лет & 0 \% ( 0 / 787 случаев) & {[} 0 ; 0.01 {]} & 0 \% (
0 / 724 случаев) & {[} 0 ; 0.01 {]} & \\
менопауза более 5 лет & 0 \% ( 0 / 787 случаев) & {[} 0 ; 0.01 {]} & 0
\% ( 0 / 724 случаев) & {[} 0 ; 0.01 {]} & \\
пременопауза & 3.304 \% ( 26 / 787 случаев) & {[} 0.02 ; 0.05 {]} &
2.762 \% ( 20 / 724 случаев) & {[} 0.02 ; 0.04 {]} & \\
оперативная менопауза & 0 \% ( 0 / 787 случаев) & {[} 0 ; 0.01 {]} & 0
\% ( 0 / 724 случаев) & {[} 0 ; 0.01 {]} & \\
\end{longtable}

\emph{Операции на молочной железе}

В группе А операций на молочной железе в анамнезе не было в 98.48\%
(775/787 случаев ) {[}95\% ДИ 0.97;0.99{]}, и было 1.52\% (12/787
случаев ) {[}95\% ДИ 0.01;0.03{]}.

В группе B операций на молочной железе в анамнезе не было в 98.34\%
(712/724 случаев ) {[}95\% ДИ 0.97;0.99{]}, в 1.66\% операции были
(12/724 случаев ) {[}95\% ДИ 0.01;0.03{]}.

\includegraphics{images/breast_surgery_before}

Разница между группами по репродуктивному статусу составила p-уровень=
0.41 . Группы распределены однородно.

\begin{longtable}[]{@{}
  >{\raggedright\arraybackslash}p{(\columnwidth - 10\tabcolsep) * \real{0.1667}}
  >{\raggedright\arraybackslash}p{(\columnwidth - 10\tabcolsep) * \real{0.1667}}
  >{\raggedright\arraybackslash}p{(\columnwidth - 10\tabcolsep) * \real{0.1667}}
  >{\raggedright\arraybackslash}p{(\columnwidth - 10\tabcolsep) * \real{0.1667}}
  >{\raggedright\arraybackslash}p{(\columnwidth - 10\tabcolsep) * \real{0.1667}}
  >{\raggedright\arraybackslash}p{(\columnwidth - 10\tabcolsep) * \real{0.1667}}@{}}
\toprule\noalign{}
\begin{minipage}[b]{\linewidth}\raggedright
Операции на МЖ
\end{minipage} & \begin{minipage}[b]{\linewidth}\raggedright
Процентная доля
\end{minipage} & \begin{minipage}[b]{\linewidth}\raggedright
95\% ДИ
\end{minipage} & \begin{minipage}[b]{\linewidth}\raggedright
Процентная доля
\end{minipage} & \begin{minipage}[b]{\linewidth}\raggedright
95\% ДИ
\end{minipage} & \begin{minipage}[b]{\linewidth}\raggedright
p-value
\end{minipage} \\
\midrule\noalign{}
\endhead
\bottomrule\noalign{}
\endlastfoot
не было операций & 98.475 \% ( 775 / 787 случаев) & {[} 0.97 ; 0.99 {]}
& 98.343 \% ( 712 / 724 случаев) & {[} 0.97 ; 0.99 {]} & \\
были операции & 1.525 \% ( 12 / 787 случаев) & {[} 0.01 ; 0.03 {]} &
1.657 \% ( 12 / 724 случаев) & {[} 0.01 ; 0.03 {]} & \\
\end{longtable}

\emph{Прием гормональных препаратов}

В группе А прием гормональных препаратов был в 20.34\% (155/787 случаев
) {[}95\% ДИ 0.17;0.23{]} и в не было в 79.66\% (607/787 случаев )
{[}95\% ДИ 0.74;0.8{]}.

В группе B прием гормональных препаратов был в 21.8\% (155/724 случаев )
{[}95\% ДИ 0.19;0.25{]} и приема препаратов не было в 78.2\% (556/724
случаев ) {[}95\% ДИ 0.74;0.8{]}.

\includegraphics{images/hormonal_medications}

\begin{longtable}[]{@{}
  >{\raggedright\arraybackslash}p{(\columnwidth - 10\tabcolsep) * \real{0.1667}}
  >{\raggedright\arraybackslash}p{(\columnwidth - 10\tabcolsep) * \real{0.1667}}
  >{\raggedright\arraybackslash}p{(\columnwidth - 10\tabcolsep) * \real{0.1667}}
  >{\raggedright\arraybackslash}p{(\columnwidth - 10\tabcolsep) * \real{0.1667}}
  >{\raggedright\arraybackslash}p{(\columnwidth - 10\tabcolsep) * \real{0.1667}}
  >{\raggedright\arraybackslash}p{(\columnwidth - 10\tabcolsep) * \real{0.1667}}@{}}
\toprule\noalign{}
\begin{minipage}[b]{\linewidth}\raggedright
Гормональные препараты
\end{minipage} & \begin{minipage}[b]{\linewidth}\raggedright
Процентная доля
\end{minipage} & \begin{minipage}[b]{\linewidth}\raggedright
95\% ДИ
\end{minipage} & \begin{minipage}[b]{\linewidth}\raggedright
Процентная доля
\end{minipage} & \begin{minipage}[b]{\linewidth}\raggedright
95\% ДИ
\end{minipage} & \begin{minipage}[b]{\linewidth}\raggedright
p-value
\end{minipage} \\
\midrule\noalign{}
\endhead
\bottomrule\noalign{}
\endlastfoot
да & 19.695 \% ( 155 / 787 случаев) & {[} 0.17 ; 0.23 {]} & 21.409 \% (
155 / 724 случаев) & {[} 0.19 ; 0.25 {]} & \\
нет & 80.305 \% ( 632 / 787 случаев) & {[} 0.77 ; 0.83 {]} & 78.591 \% (
569 / 724 случаев) & {[} 0.75 ; 0.81 {]} & \\
\end{longtable}

\emph{Генетической предрасположенность}

В группе А генетической предрасположенности не было в 68.52\% (531/787
случаев ) {[}95\% ДИ 0.64;0.71{]}, было в 31.48\% (244/787 случаев )
{[}95\% ДИ 0.28;0.34{]}.

В группе B генетическая предрасположенности не было в 68.62\% (492/724
случаев ) {[}95\% ДИ 0.64;0.71{]}, в 31.38\% была (225/724 случаев )
{[}95\% ДИ 0.28;0.35{]}.

\includegraphics{images/genetics}

\begin{longtable}[]{@{}
  >{\raggedright\arraybackslash}p{(\columnwidth - 10\tabcolsep) * \real{0.1667}}
  >{\raggedright\arraybackslash}p{(\columnwidth - 10\tabcolsep) * \real{0.1667}}
  >{\raggedright\arraybackslash}p{(\columnwidth - 10\tabcolsep) * \real{0.1667}}
  >{\raggedright\arraybackslash}p{(\columnwidth - 10\tabcolsep) * \real{0.1667}}
  >{\raggedright\arraybackslash}p{(\columnwidth - 10\tabcolsep) * \real{0.1667}}
  >{\raggedright\arraybackslash}p{(\columnwidth - 10\tabcolsep) * \real{0.1667}}@{}}
\toprule\noalign{}
\begin{minipage}[b]{\linewidth}\raggedright
Генетическая предрасположенность
\end{minipage} & \begin{minipage}[b]{\linewidth}\raggedright
Процентная доля
\end{minipage} & \begin{minipage}[b]{\linewidth}\raggedright
95\% ДИ
\end{minipage} & \begin{minipage}[b]{\linewidth}\raggedright
Процентная доля
\end{minipage} & \begin{minipage}[b]{\linewidth}\raggedright
95\% ДИ
\end{minipage} & \begin{minipage}[b]{\linewidth}\raggedright
p-value
\end{minipage} \\
\midrule\noalign{}
\endhead
\bottomrule\noalign{}
\endlastfoot
нет & 68.996 \% ( 543 / 787 случаев) & {[} 0.66 ; 0.72 {]} & 68.923 \% (
499 / 724 случаев) & {[} 0.65 ; 0.72 {]} & \\
есть & 31.004 \% ( 244 / 787 случаев) & {[} 0.28 ; 0.34 {]} & 31.077 \%
( 225 / 724 случаев) & {[} 0.28 ; 0.35 {]} & \\
\end{longtable}

\emph{Мутация BRCA}

В группе А мутация BRCA1 была выявлена в 0.76\% (6/787 случаев ) {[}95\%
ДИ 0;0.02{]}, в остальных случаях исследование не проводилось.

В группе B мутаций BRCA не выявлено было в 2.49\% (18/724 случаев )
{[}95\% ДИ 0.02;0.04{]}, в остальных случаях исследование не
проводилось.

\includegraphics{images/mutation_brca}

\begin{longtable}[]{@{}
  >{\raggedright\arraybackslash}p{(\columnwidth - 10\tabcolsep) * \real{0.1667}}
  >{\raggedright\arraybackslash}p{(\columnwidth - 10\tabcolsep) * \real{0.1667}}
  >{\raggedright\arraybackslash}p{(\columnwidth - 10\tabcolsep) * \real{0.1667}}
  >{\raggedright\arraybackslash}p{(\columnwidth - 10\tabcolsep) * \real{0.1667}}
  >{\raggedright\arraybackslash}p{(\columnwidth - 10\tabcolsep) * \real{0.1667}}
  >{\raggedright\arraybackslash}p{(\columnwidth - 10\tabcolsep) * \real{0.1667}}@{}}
\toprule\noalign{}
\begin{minipage}[b]{\linewidth}\raggedright
Мутация BRCA
\end{minipage} & \begin{minipage}[b]{\linewidth}\raggedright
Процентная доля
\end{minipage} & \begin{minipage}[b]{\linewidth}\raggedright
95\% ДИ
\end{minipage} & \begin{minipage}[b]{\linewidth}\raggedright
Процентная доля
\end{minipage} & \begin{minipage}[b]{\linewidth}\raggedright
95\% ДИ
\end{minipage} & \begin{minipage}[b]{\linewidth}\raggedright
p-value
\end{minipage} \\
\midrule\noalign{}
\endhead
\bottomrule\noalign{}
\endlastfoot
Оценка не проводилась & 99.238 \% ( 781 / 787 случаев) & {[} 0.98 ; 1
{]} & 97.514 \% ( 706 / 724 случаев) & {[} 0.96 ; 0.98 {]} & \\
BRCA1 & 0.762 \% ( 6 / 787 случаев) & {[} 0 ; 0.02 {]} & 0 \% ( 0 / 724
случаев) & {[} 0 ; 0.01 {]} & \\
мутаций не выявлено & 0 \% ( 0 / 787 случаев) & {[} 0 ; 0.01 {]} & 2.486
\% ( 18 / 724 случаев) & {[} 0.02 ; 0.04 {]} & \\
BRCA2 & 0 \% ( 0 / 787 случаев) & {[} 0 ; 0.01 {]} & 0 \% ( 0 / 724
случаев) & {[} 0 ; 0.01 {]} & \\
образование не визуал на ММГ & 0 \% ( 0 / 787 случаев) & {[} 0 ; 0.01
{]} & 0 \% ( 0 / 724 случаев) & {[} 0 ; 0.01 {]} & \\
ложноположительный результат & 0 \% ( 0 / 787 случаев) & {[} 0 ; 0.01
{]} & 0 \% ( 0 / 724 случаев) & {[} 0 ; 0.01 {]} & \\
\end{longtable}

\emph{Сторона поражения при осмотре}

В группе А сторона поражения при осмотре была правая в 7.12\% (56/787
случаев ) {[}95\% ДИ 0.05;0.09{]}, левая в 7.24\% (57/787 случаев )
{[}95\% ДИ 0.06;0.09{]}, обе стороны было 2.29\% (18/787 случаев )
{[}95\% ДИ 0.01;0.04{]}, в остальных случаях образование не
пальпировалась.

В группе B сторона поражения при осмотре была правая в 15.75\% (114/724
случаев ) {[}95\% ДИ 0.13;0.19{]}, поражение левой было в 5.8\% (42/724
случаев ) {[}95\% ДИ 0.04;0.08{]}, поражение обеих сторон обе было в
5.8\% (42/724 случаев ) {[}95\% ДИ 0.04;0.08{]}, в остальных случаях
образование не пальпировалось.

\includegraphics{images/side}

\begin{longtable}[]{@{}
  >{\raggedright\arraybackslash}p{(\columnwidth - 10\tabcolsep) * \real{0.1667}}
  >{\raggedright\arraybackslash}p{(\columnwidth - 10\tabcolsep) * \real{0.1667}}
  >{\raggedright\arraybackslash}p{(\columnwidth - 10\tabcolsep) * \real{0.1667}}
  >{\raggedright\arraybackslash}p{(\columnwidth - 10\tabcolsep) * \real{0.1667}}
  >{\raggedright\arraybackslash}p{(\columnwidth - 10\tabcolsep) * \real{0.1667}}
  >{\raggedright\arraybackslash}p{(\columnwidth - 10\tabcolsep) * \real{0.1667}}@{}}
\toprule\noalign{}
\begin{minipage}[b]{\linewidth}\raggedright
Сторона поражения
\end{minipage} & \begin{minipage}[b]{\linewidth}\raggedright
Процентная доля
\end{minipage} & \begin{minipage}[b]{\linewidth}\raggedright
95\% ДИ
\end{minipage} & \begin{minipage}[b]{\linewidth}\raggedright
Процентная доля
\end{minipage} & \begin{minipage}[b]{\linewidth}\raggedright
95\% ДИ
\end{minipage} & \begin{minipage}[b]{\linewidth}\raggedright
p-value
\end{minipage} \\
\midrule\noalign{}
\endhead
\bottomrule\noalign{}
\endlastfoot
правая & 7.116 \% ( 56 / 787 случаев) & {[} 0.05 ; 0.09 {]} & 15.746 \%
( 114 / 724 случаев) & {[} 0.13 ; 0.19 {]} & \\
левая & 7.243 \% ( 57 / 787 случаев) & {[} 0.06 ; 0.09 {]} & 5.801 \% (
42 / 724 случаев) & {[} 0.04 ; 0.08 {]} & \\
обе & 2.287 \% ( 18 / 787 случаев) & {[} 0.01 ; 0.04 {]} & 5.801 \% ( 42
/ 724 случаев) & {[} 0.04 ; 0.08 {]} & \\
не пальпируются & 83.355 \% ( 656 / 787 случаев) & {[} 0.81 ; 0.86 {]} &
72.652 \% ( 526 / 724 случаев) & {[} 0.69 ; 0.76 {]} & \\
\end{longtable}

\emph{Кожные симптомы}

В группе А и группе B во всех случаях кожных симптомов при осмотре не
было выявлено.

\emph{Симптом втягивания}

В группе А симптом втягивания соска во всех случаях был отрицательный.В
группе B симптома втягивания соска при осмотре нет было в 97.51\%
(706/724 случаев ) {[}95\% ДИ 0.96;0.98{]} и было в 2.49\% (18/724
случаев ) {[}95\% ДИ 0.02;0.04{]}.

\includegraphics{images/nipple_retraction}

\begin{longtable}[]{@{}
  >{\raggedright\arraybackslash}p{(\columnwidth - 10\tabcolsep) * \real{0.1667}}
  >{\raggedright\arraybackslash}p{(\columnwidth - 10\tabcolsep) * \real{0.1667}}
  >{\raggedright\arraybackslash}p{(\columnwidth - 10\tabcolsep) * \real{0.1667}}
  >{\raggedright\arraybackslash}p{(\columnwidth - 10\tabcolsep) * \real{0.1667}}
  >{\raggedright\arraybackslash}p{(\columnwidth - 10\tabcolsep) * \real{0.1667}}
  >{\raggedright\arraybackslash}p{(\columnwidth - 10\tabcolsep) * \real{0.1667}}@{}}
\toprule\noalign{}
\begin{minipage}[b]{\linewidth}\raggedright
Втягивание соска
\end{minipage} & \begin{minipage}[b]{\linewidth}\raggedright
Процентная доля
\end{minipage} & \begin{minipage}[b]{\linewidth}\raggedright
95\% ДИ
\end{minipage} & \begin{minipage}[b]{\linewidth}\raggedright
Процентная доля
\end{minipage} & \begin{minipage}[b]{\linewidth}\raggedright
95\% ДИ
\end{minipage} & \begin{minipage}[b]{\linewidth}\raggedright
p-value
\end{minipage} \\
\midrule\noalign{}
\endhead
\bottomrule\noalign{}
\endlastfoot
нет & 100 \% ( 787 / 787 случаев) & {[} 0.99 ; 1 {]} & 97.514 \% ( 706 /
724 случаев) & {[} 0.96 ; 0.98 {]} & \\
да & 0 \% ( 0 / 787 случаев) & {[} 0 ; 0.01 {]} & 2.486 \% ( 18 / 724
случаев) & {[} 0.02 ; 0.04 {]} & \\
\end{longtable}

\emph{Выделения из соска при осмотре}

В группе А выделений из соска при осмотре не было в 99.11\% (780/787
случаев ) {[}95\% ДИ 0.98;1{]}, кровянистые выделения были в 0.89\%
(7/787 случаев ) {[}95\% ДИ 0;0.02{]}.

В группе B выделения из соска при осмотре нет было во всех случаях.

\begin{longtable}[]{@{}
  >{\raggedright\arraybackslash}p{(\columnwidth - 12\tabcolsep) * \real{0.1429}}
  >{\raggedright\arraybackslash}p{(\columnwidth - 12\tabcolsep) * \real{0.1429}}
  >{\raggedright\arraybackslash}p{(\columnwidth - 12\tabcolsep) * \real{0.1429}}
  >{\raggedright\arraybackslash}p{(\columnwidth - 12\tabcolsep) * \real{0.1429}}
  >{\raggedright\arraybackslash}p{(\columnwidth - 12\tabcolsep) * \real{0.1429}}
  >{\raggedright\arraybackslash}p{(\columnwidth - 12\tabcolsep) * \real{0.1429}}
  >{\raggedright\arraybackslash}p{(\columnwidth - 12\tabcolsep) * \real{0.1429}}@{}}
\toprule\noalign{}
\begin{minipage}[b]{\linewidth}\raggedright
Выделения из соска
\end{minipage} & \begin{minipage}[b]{\linewidth}\raggedright
Процентная доля
\end{minipage} & \begin{minipage}[b]{\linewidth}\raggedright
95\% ДИ
\end{minipage} & \begin{minipage}[b]{\linewidth}\raggedright
Выделения из соска
\end{minipage} & \begin{minipage}[b]{\linewidth}\raggedright
Процентная доля
\end{minipage} & \begin{minipage}[b]{\linewidth}\raggedright
95\% ДИ
\end{minipage} & \begin{minipage}[b]{\linewidth}\raggedright
p-value
\end{minipage} \\
\midrule\noalign{}
\endhead
\bottomrule\noalign{}
\endlastfoot
нет & 99.111 \% ( 780 / 787 случаев) & {[} 0.98 ; 1 {]} & нет & 100 \% (
724 / 724 случаев) & {[} 0.99 ; 1 {]} & \\
кровянистые & 0.889 \% ( 7 / 787 случаев) & {[} 0 ; 0.02 {]} &
кровянистые & 0 \% ( 0 / 724 случаев) & {[} 0 ; 0.01 {]} & \\
\end{longtable}

\emph{Тип плотности}

В группе А тип плотности по ACR при осмотре был «С» в 72.27\% (555/787
случаев ) {[}95\% ДИ 0.67;0.74{]}, тип «D» в 27.73\% (213/787 случаев )
{[}95\% ДИ 0.24;0.3{]}.

В группе B тип плотности по ACR при осмотре был «C» 67.82\% (491/724
случаев ) {[}95\% ДИ 0.64;0.71{]}, тип «D» был в 32.18\% (233/724
случаев ) {[}95\% ДИ 0.29;0.36{]}.

\includegraphics{images/type_structure_acr}

\begin{longtable}[]{@{}
  >{\raggedright\arraybackslash}p{(\columnwidth - 10\tabcolsep) * \real{0.1667}}
  >{\raggedright\arraybackslash}p{(\columnwidth - 10\tabcolsep) * \real{0.1667}}
  >{\raggedright\arraybackslash}p{(\columnwidth - 10\tabcolsep) * \real{0.1667}}
  >{\raggedright\arraybackslash}p{(\columnwidth - 10\tabcolsep) * \real{0.1667}}
  >{\raggedright\arraybackslash}p{(\columnwidth - 10\tabcolsep) * \real{0.1667}}
  >{\raggedright\arraybackslash}p{(\columnwidth - 10\tabcolsep) * \real{0.1667}}@{}}
\toprule\noalign{}
\begin{minipage}[b]{\linewidth}\raggedright
Структрура ACR
\end{minipage} & \begin{minipage}[b]{\linewidth}\raggedright
Процентная доля
\end{minipage} & \begin{minipage}[b]{\linewidth}\raggedright
95\% ДИ
\end{minipage} & \begin{minipage}[b]{\linewidth}\raggedright
Процентная доля
\end{minipage} & \begin{minipage}[b]{\linewidth}\raggedright
95\% ДИ
\end{minipage} & \begin{minipage}[b]{\linewidth}\raggedright
p-value
\end{minipage} \\
\midrule\noalign{}
\endhead
\bottomrule\noalign{}
\endlastfoot
С & 72.935 \% ( 574 / 787 случаев) & {[} 0.7 ; 0.76 {]} & 67.818 \% (
491 / 724 случаев) & {[} 0.64 ; 0.71 {]} & \\
D & 27.065 \% ( 213 / 787 случаев) & {[} 0.24 ; 0.3 {]} & 32.182 \% (
233 / 724 случаев) & {[} 0.29 ; 0.36 {]} & \\
\end{longtable}

\emph{Квадрант локализации}

В группе А квадрант локализации верхне-наружный был в 4.07\% (32/787
случаев ) {[}95\% ДИ 0.03;0.06{]}, верхне-внутренний в 3.81\% (30/787
случаев ) {[}95\% ДИ 0.03;0.05{]}, нижне-наружный в 0.76\% (6/787
случаев ) {[}95\% ДИ 0;0.02{]}, нижне-внутренний в 1.52\% (12/787
случаев ) {[}95\% ДИ 0.01;0.03{]}, граница внутреннего в 1.65\% (13/787
случаев ) {[}95\% ДИ 0.01;0.03{]}, граница наружного в 0.76\% (6/787
случаев ) {[}95\% ДИ 0;0.02{]}, граница верхнего в 4.07\% (32/787
случаев ) {[}95\% ДИ 0.03;0.06{]}.

В группе B квадрант локализации верхне-наружный было 12.29\% (89/724
случаев ) {[}95\% ДИ 0.1;0.15{]}, верхне-внутренний в 4.14\% (30/724
случаев ) {[}95\% ДИ 0.03;0.06{]}, нижне-наружный в 4.14\% (30/724
случаев ) {[}95\% ДИ 0.03;0.06{]}, нижне-внутренний в 0.83\% (6/724
случаев ) {[}95\% ДИ 0;0.02{]}, граница внутреннего в 1.8\% (13/724
случаев ) {[}95\% ДИ 0.01;0.03{]}, граница наружного в 1.66\% (12/724
случаев ) {[}95\% ДИ 0.01;0.03{]}, граница верхнего в 1.66\% (12/724
случаев ) {[}95\% ДИ 0.01;0.03{]}, граница нижних в 1.8\% (13/724
случаев ) {[}95\% ДИ 0.01;0.03{]}, вся железа в 0.83\% (6/724 случаев )
{[}95\% ДИ 0;0.02{]}.

\includegraphics{images/quadrant}

\begin{longtable}[]{@{}
  >{\raggedright\arraybackslash}p{(\columnwidth - 10\tabcolsep) * \real{0.1667}}
  >{\raggedright\arraybackslash}p{(\columnwidth - 10\tabcolsep) * \real{0.1667}}
  >{\raggedright\arraybackslash}p{(\columnwidth - 10\tabcolsep) * \real{0.1667}}
  >{\raggedright\arraybackslash}p{(\columnwidth - 10\tabcolsep) * \real{0.1667}}
  >{\raggedright\arraybackslash}p{(\columnwidth - 10\tabcolsep) * \real{0.1667}}
  >{\raggedright\arraybackslash}p{(\columnwidth - 10\tabcolsep) * \real{0.1667}}@{}}
\toprule\noalign{}
\begin{minipage}[b]{\linewidth}\raggedright
Квадрант
\end{minipage} & \begin{minipage}[b]{\linewidth}\raggedright
Процентная доля
\end{minipage} & \begin{minipage}[b]{\linewidth}\raggedright
95\% ДИ
\end{minipage} & \begin{minipage}[b]{\linewidth}\raggedright
Процентная доля
\end{minipage} & \begin{minipage}[b]{\linewidth}\raggedright
95\% ДИ
\end{minipage} & \begin{minipage}[b]{\linewidth}\raggedright
p-value
\end{minipage} \\
\midrule\noalign{}
\endhead
\bottomrule\noalign{}
\endlastfoot
нет & 83.355 \% ( 656 / 787 случаев) & {[} 0.81 ; 0.86 {]} & 70.856 \% (
513 / 724 случаев) & {[} 0.67 ; 0.74 {]} & \\
верхне-наружный & 4.066 \% ( 32 / 787 случаев) & {[} 0.03 ; 0.06 {]} &
12.293 \% ( 89 / 724 случаев) & {[} 0.1 ; 0.15 {]} & \\
верхне-внутренний & 3.812 \% ( 30 / 787 случаев) & {[} 0.03 ; 0.05 {]} &
4.144 \% ( 30 / 724 случаев) & {[} 0.03 ; 0.06 {]} & \\
нижне-наружный & 0.762 \% ( 6 / 787 случаев) & {[} 0 ; 0.02 {]} & 4.144
\% ( 30 / 724 случаев) & {[} 0.03 ; 0.06 {]} & \\
нижне-внутренний & 1.525 \% ( 12 / 787 случаев) & {[} 0.01 ; 0.03 {]} &
0.829 \% ( 6 / 724 случаев) & {[} 0 ; 0.02 {]} & \\
граница внутреннего & 1.652 \% ( 13 / 787 случаев) & {[} 0.01 ; 0.03 {]}
& 1.796 \% ( 13 / 724 случаев) & {[} 0.01 ; 0.03 {]} & \\
граница наружного & 0.762 \% ( 6 / 787 случаев) & {[} 0 ; 0.02 {]} &
1.657 \% ( 12 / 724 случаев) & {[} 0.01 ; 0.03 {]} & \\
граница верхнего & 4.066 \% ( 32 / 787 случаев) & {[} 0.03 ; 0.06 {]} &
1.657 \% ( 12 / 724 случаев) & {[} 0.01 ; 0.03 {]} & \\
граница нижних & 0 \% ( 0 / 787 случаев) & {[} 0 ; 0.01 {]} & 1.796 \% (
13 / 724 случаев) & {[} 0.01 ; 0.03 {]} & \\
вся железа & 0 \% ( 0 / 787 случаев) & {[} 0 ; 0.01 {]} & 0.829 \% ( 6 /
724 случаев) & {[} 0 ; 0.02 {]} & \\
1/2 железы & 0 \% ( 0 / 787 случаев) & {[} 0 ; 0.01 {]} & 0 \% ( 0 / 724
случаев) & {[} 0 ; 0.01 {]} & \\
центральный квадрант & 0 \% ( 0 / 787 случаев) & {[} 0 ; 0.01 {]} & 0 \%
( 0 / 724 случаев) & {[} 0 ; 0.01 {]} & \\
\end{longtable}

\hypertarget{ux43eux431ux449ux435ux435-ux43eux43fux438ux441ux430ux43dux438ux435-ux432ux44bux431ux43eux440ux43aux438-ux43fux430ux446ux438ux435ux43dux442ux43eux43a-40-ux43bux435ux442-ux438-ux441ux442ux430ux440ux448ux435}{%
\subsubsection{2.5.3 Общее описание выборки пациенток 40 лет и
старше}\label{ux43eux431ux449ux435ux435-ux43eux43fux438ux441ux430ux43dux438ux435-ux432ux44bux431ux43eux440ux43aux438-ux43fux430ux446ux438ux435ux43dux442ux43eux43a-40-ux43bux435ux442-ux438-ux441ux442ux430ux440ux448ux435}}

С февраля 2019 по май 2023 года было исследовано 1283 пациенток с
возрастом 40 лет и старше.

Медиана возраста пациенток выборке после 40 лет составил 49 {[}Q1-Q3:
45;56{]} лет. Минимальный возраст составил 40 лет. Максимальный возраст
составил 79 лет. Медиана роста пациенток выборке после 40 лет составил
166 {[}Q1-Q3: 164;168{]} см. Медиана веса пациенток выборке после 40 лет
составил 65 {[}Q1-Q3: 60;73.25{]} кг. В выборке после 40 лет были
поставлены диагнозы Без патологии было 5.61\% (72/1283 случаев) {[}95\%
ДИ 0.04;0.07{]}.

В выборке 40 лет и старше были поставлены диагнозы при первичном
осмотре: «Рак» было 3.51\% (45/1283 случаев) {[}95\% ДИ 0.03;0.05{]},
«Фиброаденома» в 10.05\% (129/1283 случаев) {[}95\% ДИ 0.08;0.12{]},
«Фиброзно-кистозная болезнь» - 31.57\% (405/1283 случаев) {[}95\% ДИ
0.29;0.34{]}, «Киста» было 2.96\% (38/1283 случаев) {[}95\% ДИ
0.02;0.04{]}, «Листовидная опухоль» было 1.71\% (22/1283 случаев)
{[}95\% ДИ 0.01;0.03{]}, ``Диффузный фиброаденоматоз'' - 41.86\%
(537/1283 случаев) {[}95\% ДИ 0.39;0.45{]}, «Локализованный
фиброаденоматоз» был в 2.49\% (32/1283 случаев) {[}95\% ДИ 0.02;0.04{]},
«Мастит» был в 0.23\% (3/1283 случаев) {[}95\% ДИ 0;0.01{]}.

Были выявлены следующие жалобы: выделения из соска было в 0.47\% (6/1283
случаев) {[}95\% ДИ 0;0.01{]}, боль была в 4.13\% (53/1283 случаев)
{[}95\% ДИ 0.03;0.05{]}, уплотнение было в 14.34\% (184/1283 случаев)
{[}95\% ДИ 0.12;0.16{]}, жалоб не было в 81.06\% (1040/1283 случаев)
{[}95\% ДИ 0.79;0.83{]}.

Репродуктивный возраст был в 22.37\% (287/1283 случаев) {[}95\% ДИ
0.2;0.25{]}, менопауза до 5 лет была в 18.47\% (237/1283 случаев)
{[}95\% ДИ 0.16;0.21{]}, менопауза более 5 лет была в 24.94\% (320/1283
случаев) {[}95\% ДИ 0.23;0.27{]}, пременопауза - 33.9\% (435/1283
случаев) {[}95\% ДИ 0.31;0.37{]}, оперативная менопауза - 0.31\% (4/1283
случаев) {[}95\% ДИ 0;0.01{]}.

В выборке после 40 лет операций на молочной железе в анамнезе не было в
99.14\% (1272/1283 случаев) {[}95\% ДИ 0.98;1{]}, было в 0.86\% (11/1283
случаев) {[}95\% ДИ 0;0.02{]}.

В выборке после 40 лет прием гормональных препаратов был в 20.83\%
(265/1283 случаев) {[}95\% ДИ 0.18;0.23{]}, не было 79.17\% (1007/1283
случаев) {[}95\% ДИ 0.76;0.81{]}.

В выборке после 40 лет генетической предрасположенности не было в 63.2\%
(802/1283 случаев) {[}95\% ДИ 0.6;0.65{]}, была в 36.8\% (467/1283
случаев) {[}95\% ДИ 0.34;0.39{]}.

При оценке мутация BRCA ее не выявлено было в 1.79\% (23/1283 случаев)
{[}95\% ДИ 0.01;0.03{]}, мутация BRCA2 была в 0.31\% (4/1283 случаев)
{[}95\% ДИ 0;0.01{]}, образование не визуализировалось на ММГ в 0.86\%
(11/1283 случаев) {[}95\% ДИ 0;0.02{]}, ложноположительный результат
мутации BRCA был в 0.31\% (4/1283 случаев) {[}95\% ДИ 0;0.01{]}.

В выборке после 40 лет сторона поражения при осмотре была правая в 9.9\%
(127/1283 случаев) {[}95\% ДИ 0.08;0.12{]}, левая - 8.96\% (115/1283
случаев) {[}95\% ДИ 0.07;0.11{]}, поражение обеих сторон - 4.75\%
(61/1283 случаев) {[}95\% ДИ 0.04;0.06{]}, в остальных случаях
образование не пальпировалось - 76.38\% (980/1283 случаев) {[}95\% ДИ
0.74;0.79{]}.

Кожных симптомы при осмотре нет было в 99.06\% (1271/1283 случаев)
{[}95\% ДИ 0.98;0.99{]}, отек был в 0.23\% (3/1283 случаев) {[}95\% ДИ
0;0.01{]}, гиперемия была в 0.7\% (9/1283 случаев) {[}95\% ДИ 0;0.01{]}.

Втягивание соска при осмотре нет было в 96.02\% (1232/1283 случаев)
{[}95\% ДИ 0.95;0.97{]} и было в 3.98\% (51/1283 случаев) {[}95\% ДИ
0.03;0.05{]}.

Выделения из соска при осмотре не было в 98.91\% (1269/1283 случаев)
{[}95\% ДИ 0.98;0.99{]}, светлые выделения были в 0.62\% (8/1283
случаев) {[}95\% ДИ 0;0.01{]}, кровянистые были в 0.47\% (6/1283
случаев) {[}95\% ДИ 0;0.01{]}.

Тип плотности по ACR при осмотре был «С» в 87.96\% (1125/1283 случаев)
{[}95\% ДИ 0.86;0.89{]}, тип «D» былл в 12.04\% (154/1283 случаев)
{[}95\% ДИ 0.1;0.14{]}.

Квадрант локализации верхне-наружный был в 10.52\% (135/1283 случаев)
{[}95\% ДИ 0.09;0.12{]}, верхне-внутренний в 4.68\% (60/1283 случаев)
{[}95\% ДИ 0.04;0.06{]}, нижне-наружный в 1.95\% (25/1283 случаев)
{[}95\% ДИ 0.01;0.03{]}, нижне-внутренний в 1.48\% (19/1283 случаев)
{[}95\% ДИ 0.01;0.02{]}, граница наружного была в 2.18\% (28/1283
случаев) {[}95\% ДИ 0.01;0.03{]}, граница верхнего в 1.64\% (21/1283
случаев) {[}95\% ДИ 0.01;0.03{]}, граница нижних было 0.78\% (10/1283
случаев) {[}95\% ДИ 0;0.01{]}, вся железа в 0.55\% (7/1283 случаев)
{[}95\% ДИ 0;0.01{]}, центральный квадрант в 0.55\% (7/1283 случаев)
{[}95\% ДИ 0;0.01{]}.

\hypertarget{ux43eux446ux435ux43dux43aux430-ux43eux434ux43dux43eux440ux43eux434ux43dux43eux441ux442ux438-ux43fux430ux446ux438ux435ux43dux442ux43eux43a-40-ux43bux435ux442-ux438-ux441ux442ux430ux440ux448ux435}{%
\subsubsection{2.5.4 Оценка однородности пациенток 40 лет и
старше}\label{ux43eux446ux435ux43dux43aux430-ux43eux434ux43dux43eux440ux43eux434ux43dux43eux441ux442ux438-ux43fux430ux446ux438ux435ux43dux442ux43eux43a-40-ux43bux435ux442-ux438-ux441ux442ux430ux440ux448ux435}}

Пациенток в группе C в исследование вошло 628 пациентов. Медиана
возраста пациенток в группе C составила 49.5 {[}Q1-Q3: 45;56{]} лет.
Минимальный возраст составил 40 лет. Максимальный возраст составил 79
лет. Медиана роста составила 167 {[}Q1-Q3: 164;169{]} см. Медиана веса
составила 65 {[}Q1-Q3: 60;72{]} кг. Пациенток в группу D вошло 655
пациентов. Медиана возраста составила 49 {[}Q1-Q3: 45;56.5{]} лет.
Минимальный возраст составил 40 лет. Максимальный возраст составил 72
лет. Медиана роста составила 165 {[}Q1-Q3: 164;168{]} см. Медиана веса
составила 65 {[}Q1-Q3: 59;75{]} кг.

\emph{Первичный диагноз}

В группе C были поставлены следующие диагнозы: «Без патологии» было
5.1\% (32/628 случаев ) {[}95\% ДИ 0.04;0.07{]}, «Рак» было 1.59\%
(10/628 случаев ) {[}95\% ДИ 0.01;0.03{]}, «Фиброаденома» в 7.96\%
(50/628 случаев ) {[}95\% ДИ 0.06;0.1{]}, «Фиброзно-кистозная болезнь» в
34.87\% (219/628 случаев ) {[}95\% ДИ 0.31;0.39{]}, «Киста» в 2.71\%
(17/628 случаев ) {[}95\% ДИ 0.02;0.04{]}, «Листовидная опухоль» в
1.11\% (7/628 случаев ) {[}95\% ДИ 0;0.02{]}, «Диффузный
фиброаденоматоз» в 44.43\% (279/628 случаев ) {[}95\% ДИ 0.41;0.48{]},
«Локализованный фиброаденоматоз» в 2.23\% (14/628 случаев ) {[}95\% ДИ
0.01;0.04{]}. В группе D были поставлены диагнозы: Без патологии было
6.11\% (40/655 случаев ) {[}95\% ДИ 0.04;0.08{]}, «Рак» было 5.34\%
(35/655 случаев ) {[}95\% ДИ 0.04;0.07{]}, «Фиброаденома'' в 12.06\%
(79/655 случаев ) {[}95\% ДИ 0.1;0.15{]}, «Фиброзно-кистозная болезнь» в
28.4\% (186/655 случаев ) {[}95\% ДИ 0.25;0.32{]}, «Киста» в 3.21\%
(21/655 случаев ) {[}95\% ДИ 0.02;0.05{]}, «Листовидная опухоль» в
2.29\% (15/655 случаев ) {[}95\% ДИ 0.01;0.04{]}, «Диффузный
фиброаденоматоз» в 39.39\% (258/655 случаев ) {[}95\% ДИ 0.36;0.43{]},
«Локализованный фиброаденоматоз» в 2.75\% (18/655 случаев ) {[}95\% ДИ
0.02;0.04{]},«Мастит» в 0.46\% (3/655 случаев ) {[}95\% ДИ 0;0.01{]}.

\includegraphics{images/diagnosis_primary}

\emph{Жалобы при осмотре}

В группе C были выявлены следующие жалобы нет было 88.38\% (555/628
случаев ) {[}95\% ДИ 0.86;0.91{]}, выделения из соска были в 0.96\%
(6/628 случаев ) {[}95\% ДИ 0;0.02{]}, боль была в 4.94\% (31/628
случаев ) {[}95\% ДИ 0.03;0.07{]} и уплотнение было в 5.73\% (36/628
случаев ) {[}95\% ДИ 0.04;0.08{]}. В группе D при осмотре жалоб не было
в 74.05\% (485/655 случаев ) {[}95\% ДИ 0.7;0.77{]}, жалобы на боль была
в 3.36\% (22/655 случаев ) {[}95\% ДИ 0.02;0.05{]}, на уплотнение в
22.6\% (148/655 случаев ) {[}95\% ДИ 0.19;0.26{]}.

\includegraphics{images/complaints-01}

\emph{Репродуктивный статус}

В группе C репродуктивный возраст был в 21.5\% (135/628 случаев )
{[}95\% ДИ 0.18;0.25{]}, менопауза до 5 лет в 21.97\% (138/628 случаев )
{[}95\% ДИ 0.19;0.25{]}, менопауза более 5 лет в 23.89\% (150/628
случаев ) {[}95\% ДИ 0.21;0.27{]}, пременопауза в 32.01\% (201/628
случаев ) {[}95\% ДИ 0.28;0.36{]}, оперативная менопауза в 0.64\% (4/628
случаев ) {[}95\% ДИ 0;0.02{]}. В группе D репродуктивный возраст был в
23.21\% (152/655 случаев ) {[}95\% ДИ 0.2;0.27{]}, менопауза до 5 лет в
15.11\% (99/655 случаев ) {[}95\% ДИ 0.13;0.18{]}, менопауза более 5 лет
в 25.95\% (170/655 случаев ) {[}95\% ДИ 0.23;0.3{]}, пременопауза в
35.73\% (234/655 случаев ) {[}95\% ДИ 0.32;0.4{]}.

\includegraphics{images/satus_reproductive-01}

\emph{Операции на молочной железе}

В группе C операций на молочной железе в анамнезе не было в 98.25\%
(617/628 случаев ) {[}95\% ДИ 0.97;0.99{]} и в 1.75\% было (11/628
случаев ) {[}95\% ДИ 0.01;0.03{]}. В группе D операций на молочной
железе в анамнезе не было операций было во всех случаях.

\includegraphics{images/breast_surgery_before-02}

\emph{Гормональные препараты}

В группе C прием гормональных препаратов был в 22.76\% (142/628 случаев
) {[}95\% ДИ 0.19;0.26{]}, в 77.24\% не было (482/628 случаев ) {[}95\%
ДИ 0.73;0.8{]}. В группе D прием гормональных препаратов был в 18.98\%
(123/655 случаев ) {[}95\% ДИ 0.16;0.22{]}, и не было в 81.02\% (525/655
случаев ) {[}95\% ДИ 0.77;0.83{]}.

\includegraphics{images/hormonal_medications-01}

\emph{Генетическая предрасположенность}

В группе C генетическая предрасположенности не было в 59.71\% (369/628
случаев ) {[}95\% ДИ 0.55;0.63{]}, в 40.29\% было (249/628 случаев )
{[}95\% ДИ 0.36;0.44{]}. В группе D генетической предрасположенности не
было в 66.51\% (433/655 случаев ) {[}95\% ДИ 0.62;0.7{]} и в 33.49\%
было (218/655 случаев ) {[}95\% ДИ 0.3;0.37{]}.

\includegraphics{images/genetics-02}

\emph{Мутации BRCA}

В группе C было при исследовании мутации BRCA ее не выявлено было 1.27\%
(8/628 случаев ) {[}95\% ДИ 0.01;0.03{]}, в остальных случаях
исследование не проводилось. В группе D при исследовании мутации BRCA не
было выявлено 2.29\% (15/655 случаев ) {[}95\% ДИ 0.01;0.04{]}, мутация
BRCA2 в 0.61\% (4/655 случаев ) {[}95\% ДИ 0;0.02{]}, образование не
визуализировалось на ММГ было 1.68\% (11/655 случаев ) {[}95\% ДИ
0.01;0.03{]}, ложноположительный результат был в 0.61\% (4/655 случаев )
{[}95\% ДИ 0;0.02{]}.

\includegraphics{images/mutation_brca-01}

\emph{Сторона поражения}

В группе C сторона поражения при осмотре правая была в 6.69\% (42/628
случаев ) {[}95\% ДИ 0.05;0.09{]}, левая было 3.66\% (23/628 случаев )
{[}95\% ДИ 0.02;0.06{]}, поражение обеих сторон было 3.18\% (20/628
случаев ) {[}95\% ДИ 0.02;0.05{]}, в остальных случаях образование не
пальпировалось. В группе D сторона поражения при осмотре была правая в
12.98\% (85/655 случаев ) {[}95\% ДИ 0.11;0.16{]}, левая в 14.05\%
(92/655 случаев ) {[}95\% ДИ 0.12;0.17{]}, поражение обеих сторон было в
6.26\% (41/655 случаев ) {[}95\% ДИ 0.05;0.08{]}, в остальных случаях
образование не пальпировалось.

\includegraphics{images/side-01}

\emph{Кожные симптомы}

В группе C кожных симптомов при осмотре не было в 99.52\% (625/628
случаев ) {[}95\% ДИ 0.98;1{]}, гиперемия была в 0.48\% (3/628 случаев )
{[}95\% ДИ 0;0.02{]}. В группе D кожных симптомов при осмотре не было в
98.63\% (646/655 случаев ) {[}95\% ДИ 0.97;0.99{]}, отек был 0.46\%
(3/655 случаев ) {[}95\% ДИ 0;0.01{]}, гиперемия была в 0.92\% (6/655
случаев ) {[}95\% ДИ 0;0.02{]}.

\includegraphics{images/skin_symptoms}

\emph{Симптом втягивания соска}

В группе C симптома втягивания соска при осмотре не было в 95.54\%
(600/628 случаев ) {[}95\% ДИ 0.94;0.97{]}, и было в 4.46\% (28/628
случаев ) {[}95\% ДИ 0.03;0.06{]}. В группе D симптома втягивания соска
при осмотре нет было в 96.49\% (632/655 случаев ) {[}95\% ДИ
0.95;0.98{]}, симптом был в 3.51\% (23/655 случаев ) {[}95\% ДИ
0.02;0.05{]}.

\includegraphics{images/nipple_retraction-01}

\emph{Выделения из соска}

В группе C выделений из соска при осмотре не было в 99.52\% (625/628
случаев ) {[}95\% ДИ 0.98;1{]}, кровянистые были в 0.48\% (3/628 случаев
) {[}95\% ДИ 0;0.02{]}. В группе D выделения из соска при осмотре не был
в 98.32\% (644/655 случаев ) {[}95\% ДИ 0.97;0.99{]}, светлые выделения
не были 1.22\% (8/655 случаев ) {[}95\% ДИ 0.01;0.02{]}, кровянистые
были в 0.46\% (3/655 случаев ) {[}95\% ДИ 0;0.01{]}.

\includegraphics{images/nipple_release}

\emph{Тип плотности}

В группе C тип плотности по ACR при осмотре был «C» 87.34\% (545/628
случаев ) {[}95\% ДИ 0.84;0.89{]}, тип «D» было 12.66\% (79/628 случаев
) {[}95\% ДИ 0.1;0.15{]}. В группе D тип плотности по ACR при осмотре
был «С» 88.55\% (580/655 случаев ) {[}95\% ДИ 0.86;0.91{]}, тип «D» былл
11.45\% (75/655 случаев ) {[}95\% ДИ 0.09;0.14{]}.

\includegraphics{images/type_structure_acr-01}

\emph{Квадрант локализации}

В группе C квадрант локализации верхне-наружный было 7.48\% (47/628
случаев ) {[}95\% ДИ 0.06;0.1{]}, верхне-внутренний в 1.75\% (11/628
случаев ) {[}95\% ДИ 0.01;0.03{]}, нижне-наружный в 1.11\% (7/628
случаев ) {[}95\% ДИ 0;0.02{]}, нижне-внутренний в 1.91\% (12/628
случаев ) {[}95\% ДИ 0.01;0.03{]}, граница наружного в 0.96\% (6/628
случаев ) {[}95\% ДИ 0;0.02{]},граница верхнего в 1.75\% (11/628 случаев
) {[}95\% ДИ 0.01;0.03{]}, граница нижних в 0.48\% (3/628 случаев )
{[}95\% ДИ 0;0.02{]}. В группе D квадрант локализации верхне-наружный
было в13.44\% (88/655 случаев ) {[}95\% ДИ 0.11;0.16{]},
верхне-внутренний в 7.48\% (49/655 случаев ) {[}95\% ДИ 0.06;0.1{]},
нижне-наружный в 2.75\% (18/655 случаев ) {[}95\% ДИ 0.02;0.04{]},
нижне-внутренний в 1.07\% (7/655 случаев ) {[}95\% ДИ 0;0.02{]}, граница
наружного в 3.36\% (22/655 случаев ) {[}95\% ДИ 0.02;0.05{]}, граница
верхнего в 1.53\% (10/655 случаев ) {[}95\% ДИ 0.01;0.03{]}, граница
нижних в 1.07\% (7/655 случаев ) {[}95\% ДИ 0;0.02{]}, вся железа в
1.07\% (7/655 случаев ) {[}95\% ДИ 0;0.02{]}., центральный квадрант было
1.07\% (7/655 случаев ) {[}95\% ДИ 0;0.02{]}.

\includegraphics{images/quadrant-01}

\hypertarget{ux43eux43fux438ux441ux430ux43dux438ux435-ux433ux440ux443ux43fux43fux44b-a}{%
\subsubsection{2.5.3 Описание группы
A}\label{ux43eux43fux438ux441ux430ux43dux438ux435-ux433ux440ux443ux43fux43fux44b-a}}

В группе А операций на молочной железе в анамнезе не было в 98.48\%
(775/787 случаев ) {[}95\% ДИ 0.97;0.99{]}, и было 1.52\% (12/787
случаев ) {[}95\% ДИ 0.01;0.03{]}.

В группе А прием гормональных препаратов был в 20.34\% (155/787 случаев
) {[}95\% ДИ 0.17;0.23{]} и в не было в 79.66\% (607/787 случаев )
{[}95\% ДИ 0.74;0.8{]}.

В группе А генетической предрасположенности не было в 68.52\% (531/787
случаев ) {[}95\% ДИ 0.64;0.71{]}, было в 31.48\% (244/787 случаев )
{[}95\% ДИ 0.28;0.34{]}.

В группе А мутация BRCA1 была выявлена в 0.76\% (6/787 случаев ) {[}95\%
ДИ 0;0.02{]}, в остальных случаях исследование не проводилось.

В группе А сторона поражения при осмотре была правая в 7.12\% (56/787
случаев ) {[}95\% ДИ 0.05;0.09{]}, левая в 7.24\% (57/787 случаев )
{[}95\% ДИ 0.06;0.09{]}, обе стороны было 2.29\% (18/787 случаев )
{[}95\% ДИ 0.01;0.04{]}, в остальных случаях образование не
пальпировалась.

В группе А во всех случаях кожных симптомом при осмотре не было
выявлено.

В группе А симптом втягивания соска во всех случаях был отрицательный

В группе А выделений из соска при осмотре не было в 99.11\% (780/787
случаев ) {[}95\% ДИ 0.98;1{]}, кровянистые выделения были в 0.89\%
(7/787 случаев ) {[}95\% ДИ 0;0.02{]}.

В группе А тип плотности по ACR при осмотре был «С» в 72.27\% (555/787
случаев ) {[}95\% ДИ 0.67;0.74{]}, тип «D» в 27.73\% (213/787 случаев )
{[}95\% ДИ 0.24;0.3{]}.

В группе А квадрант локализации верхне-наружный был в 4.07\% (32/787
случаев ) {[}95\% ДИ 0.03;0.06{]}, верхне-внутренний в 3.81\% (30/787
случаев ) {[}95\% ДИ 0.03;0.05{]}, нижне-наружный в 0.76\% (6/787
случаев ) {[}95\% ДИ 0;0.02{]}, нижне-внутренний в 1.52\% (12/787
случаев ) {[}95\% ДИ 0.01;0.03{]}, граница внутреннего в 1.65\% (13/787
случаев ) {[}95\% ДИ 0.01;0.03{]}, граница наружного в 0.76\% (6/787
случаев ) {[}95\% ДИ 0;0.02{]}, граница верхнего в 4.07\% (32/787
случаев ) {[}95\% ДИ 0.03;0.06{]}.

\hypertarget{ux43eux43fux438ux441ux430ux43dux438ux435-ux433ux440ux443ux43fux43fux44b-b}{%
\subsubsection{2.5.4 Описание группы
B}\label{ux43eux43fux438ux441ux430ux43dux438ux435-ux433ux440ux443ux43fux43fux44b-b}}

В группе B операций на молочной железе в анамнезе не было в 98.34\%
(712/724 случаев ) {[}95\% ДИ 0.97;0.99{]}, в 1.66\% операции были
(12/724 случаев ) {[}95\% ДИ 0.01;0.03{]}.

В группе B прием гормональных препаратов был в 21.8\% (155/724 случаев )
{[}95\% ДИ 0.19;0.25{]} и приема препаратов не было в 78.2\% (556/724
случаев ) {[}95\% ДИ 0.74;0.8{]}.

В группе B генетическая предрасположенности не было в 68.62\% (492/724
случаев ) {[}95\% ДИ 0.64;0.71{]}, в 31.38\% была (225/724 случаев )
{[}95\% ДИ 0.28;0.35{]}.

В группе B мутаций BRCA не выявлено было в 2.49\% (18/724 случаев )
{[}95\% ДИ 0.02;0.04{]}, в остальных случаях исследование не
проводилось.

В группе B сторона поражения при осмотре была правая в 15.75\% (114/724
случаев ) {[}95\% ДИ 0.13;0.19{]}, поражение левой было в 5.8\% (42/724
случаев ) {[}95\% ДИ 0.04;0.08{]}, поражение обеих сторон обе было в
5.8\% (42/724 случаев ) {[}95\% ДИ 0.04;0.08{]}, в остальных случаях
образование не пальпировалось.

В группе B кожные симптомы при осмотре нет было выявлено во всех
случаях.

В группе B симптома втягивания соска при осмотре нет было в 97.51\%
(706/724 случаев ) {[}95\% ДИ 0.96;0.98{]} и было в 2.49\% (18/724
случаев ) {[}95\% ДИ 0.02;0.04{]}.

В группе B выделения из соска при осмотре нет было во всех случаях.

В группе B тип плотности по ACR при осмотре был «C» 67.82\% (491/724
случаев ) {[}95\% ДИ 0.64;0.71{]}, тип «D» был в 32.18\% (233/724
случаев ) {[}95\% ДИ 0.29;0.36{]}.

В группе B квадрант локализации верхне-наружный было 12.29\% (89/724
случаев ) {[}95\% ДИ 0.1;0.15{]}, верхне-внутренний в 4.14\% (30/724
случаев ) {[}95\% ДИ 0.03;0.06{]}, нижне-наружный в 4.14\% (30/724
случаев ) {[}95\% ДИ 0.03;0.06{]}, нижне-внутренний в 0.83\% (6/724
случаев ) {[}95\% ДИ 0;0.02{]}, граница внутреннего в 1.8\% (13/724
случаев ) {[}95\% ДИ 0.01;0.03{]}, граница наружного в 1.66\% (12/724
случаев ) {[}95\% ДИ 0.01;0.03{]}, граница верхнего в 1.66\% (12/724
случаев ) {[}95\% ДИ 0.01;0.03{]}, граница нижних в 1.8\% (13/724
случаев ) {[}95\% ДИ 0.01;0.03{]}, вся железа в 0.83\% (6/724 случаев )
{[}95\% ДИ 0;0.02{]}.

\hypertarget{ux43eux43fux438ux441ux430ux43dux438ux435-ux433ux440ux443ux43fux43fux44b-c}{%
\subsubsection{2.5.5 Описание группы
C}\label{ux43eux43fux438ux441ux430ux43dux438ux435-ux433ux440ux443ux43fux43fux44b-c}}

Пациенток в группе C в исследование вошло 628 пациентов. Медиана
возраста пациенток в группе C составила 49.5 {[}Q1-Q3: 45;56{]} лет.
Минимальный возраст составил 40 лет. Максимальный возраст составил 79
лет.

Медиана роста составила 167 {[}Q1-Q3: 164;169{]} см. Медиана веса
составила 65 {[}Q1-Q3: 60;72{]} кг.

В группе C были поставлены следующие диагнозы: «Без патологии» было
5.1\% (32/628 случаев ) {[}95\% ДИ 0.04;0.07{]}, «Рак» было 1.59\%
(10/628 случаев ) {[}95\% ДИ 0.01;0.03{]}, «Фиброаденома» в 7.96\%
(50/628 случаев ) {[}95\% ДИ 0.06;0.1{]}, «Фиброзно-кистозная болезнь» в
34.87\% (219/628 случаев ) {[}95\% ДИ 0.31;0.39{]}, «Киста» в 2.71\%
(17/628 случаев ) {[}95\% ДИ 0.02;0.04{]}, «Листовидная опухоль» в
1.11\% (7/628 случаев ) {[}95\% ДИ 0;0.02{]}, «Диффузный
фиброаденоматоз» в 44.43\% (279/628 случаев ) {[}95\% ДИ 0.41;0.48{]},
«Локализованный фиброаденоматоз» в 2.23\% (14/628 случаев ) {[}95\% ДИ
0.01;0.04{]}.

В группе C были выявлены следующие жалобы нет было 88.38\% (555/628
случаев ) {[}95\% ДИ 0.86;0.91{]}, выделения из соска были в 0.96\%
(6/628 случаев ) {[}95\% ДИ 0;0.02{]}, боль была в 4.94\% (31/628
случаев ) {[}95\% ДИ 0.03;0.07{]} и уплотнение было в 5.73\% (36/628
случаев ) {[}95\% ДИ 0.04;0.08{]}.

В группе C репродуктивный возраст был в 21.5\% (135/628 случаев )
{[}95\% ДИ 0.18;0.25{]}, менопауза до 5 лет в 21.97\% (138/628 случаев )
{[}95\% ДИ 0.19;0.25{]}, менопауза более 5 лет в 23.89\% (150/628
случаев ) {[}95\% ДИ 0.21;0.27{]}, пременопауза в 32.01\% (201/628
случаев ) {[}95\% ДИ 0.28;0.36{]}, оперативная менопауза в 0.64\% (4/628
случаев ) {[}95\% ДИ 0;0.02{]}.

В группе C операций на молочной железе в анамнезе не было в 98.25\%
(617/628 случаев ) {[}95\% ДИ 0.97;0.99{]} и в 1.75\% было (11/628
случаев ) {[}95\% ДИ 0.01;0.03{]}.

В группе C прием гормональных препаратов был в 22.76\% (142/628 случаев
) {[}95\% ДИ 0.19;0.26{]}, в 77.24\% не было (482/628 случаев ) {[}95\%
ДИ 0.73;0.8{]}.

В группе C генетическая предрасположенности не было в 59.71\% (369/628
случаев ) {[}95\% ДИ 0.55;0.63{]}, в 40.29\% было (249/628 случаев )
{[}95\% ДИ 0.36;0.44{]}.

В группе C было при исследовании мутации BRCA ее не выявлено было 1.27\%
(8/628 случаев ) {[}95\% ДИ 0.01;0.03{]}, в остальных случаях
исследование не проводилось.

В группе C сторона поражения при осмотре правая была в 6.69\% (42/628
случаев ) {[}95\% ДИ 0.05;0.09{]}, левая было 3.66\% (23/628 случаев )
{[}95\% ДИ 0.02;0.06{]}, поражение обеих сторон было 3.18\% (20/628
случаев ) {[}95\% ДИ 0.02;0.05{]}, в остальных случаях образование не
пальпировалось.

В группе C кожных симптомов при осмотре не было в 99.52\% (625/628
случаев ) {[}95\% ДИ 0.98;1{]}, гиперемия была в 0.48\% (3/628 случаев )
{[}95\% ДИ 0;0.02{]}.

В группе C симптома втягивания соска при осмотре не было в 95.54\%
(600/628 случаев ) {[}95\% ДИ 0.94;0.97{]}, и было в 4.46\% (28/628
случаев ) {[}95\% ДИ 0.03;0.06{]}.

В группе C выделений из соска при осмотре не было в 99.52\% (625/628
случаев ) {[}95\% ДИ 0.98;1{]}, кровянистые были в 0.48\% (3/628 случаев
) {[}95\% ДИ 0;0.02{]}.

В группе C тип плотности по ACR при осмотре был «C» 87.34\% (545/628
случаев ) {[}95\% ДИ 0.84;0.89{]}, тип «D» было 12.66\% (79/628 случаев
) {[}95\% ДИ 0.1;0.15{]}.

В группе C квадрант локализации верхне-наружный было 7.48\% (47/628
случаев ) {[}95\% ДИ 0.06;0.1{]}, верхне-внутренний в 1.75\% (11/628
случаев ) {[}95\% ДИ 0.01;0.03{]}, нижне-наружный в 1.11\% (7/628
случаев ) {[}95\% ДИ 0;0.02{]}, нижне-внутренний в 1.91\% (12/628
случаев ) {[}95\% ДИ 0.01;0.03{]}, граница наружного в 0.96\% (6/628
случаев ) {[}95\% ДИ 0;0.02{]},граница верхнего в 1.75\% (11/628 случаев
) {[}95\% ДИ 0.01;0.03{]}, граница нижних в 0.48\% (3/628 случаев )
{[}95\% ДИ 0;0.02{]}.

\hypertarget{ux43eux43fux438ux441ux430ux43dux438ux435-ux433ux440ux443ux43fux43fux44b-d}{%
\subsubsection{2.5.6 Описание группы
D}\label{ux43eux43fux438ux441ux430ux43dux438ux435-ux433ux440ux443ux43fux43fux44b-d}}

Пациенток в группу D вошло 655 пациентов. Медиана возраста составила 49
{[}Q1-Q3: 45;56.5{]} лет. Минимальный возраст составил 40 лет.
Максимальный возраст составил 72 лет. Медиана роста составила 165
{[}Q1-Q3: 164;168{]} см. Медиана веса составила 65 {[}Q1-Q3: 59;75{]}
кг.

В группе D были поставлены диагнозы: Без патологии было 6.11\% (40/655
случаев ) {[}95\% ДИ 0.04;0.08{]}, «Рак» было 5.34\% (35/655 случаев )
{[}95\% ДИ 0.04;0.07{]}, «Фиброаденома'' в 12.06\% (79/655 случаев )
{[}95\% ДИ 0.1;0.15{]}, «Фиброзно-кистозная болезнь» в 28.4\% (186/655
случаев ) {[}95\% ДИ 0.25;0.32{]}, «Киста» в 3.21\% (21/655 случаев )
{[}95\% ДИ 0.02;0.05{]}, «Листовидная опухоль» в 2.29\% (15/655 случаев
) {[}95\% ДИ 0.01;0.04{]}, «Диффузный фиброаденоматоз» в 39.39\%
(258/655 случаев ) {[}95\% ДИ 0.36;0.43{]}, «Локализованный
фиброаденоматоз» в 2.75\% (18/655 случаев ) {[}95\% ДИ
0.02;0.04{]},«Мастит» в 0.46\% (3/655 случаев ) {[}95\% ДИ 0;0.01{]}.

В группе D при осмотре жалоб не было в 74.05\% (485/655 случаев )
{[}95\% ДИ 0.7;0.77{]}, жалобы на боль была в 3.36\% (22/655 случаев )
{[}95\% ДИ 0.02;0.05{]}, на уплотнение в 22.6\% (148/655 случаев )
{[}95\% ДИ 0.19;0.26{]}.

В группе D репродуктивный возраст был в 23.21\% (152/655 случаев )
{[}95\% ДИ 0.2;0.27{]}, менопауза до 5 лет в 15.11\% (99/655 случаев )
{[}95\% ДИ 0.13;0.18{]}, менопауза более 5 лет в 25.95\% (170/655
случаев ) {[}95\% ДИ 0.23;0.3{]}, пременопауза в 35.73\% (234/655
случаев ) {[}95\% ДИ 0.32;0.4{]}.

В группе D операций на молочной железе в анамнезе не было операций было
во всех случаях.

В группе D прием гормональных препаратов был в 18.98\% (123/655 случаев
) {[}95\% ДИ 0.16;0.22{]}, и не было в 81.02\% (525/655 случаев )
{[}95\% ДИ 0.77;0.83{]}.

В группе D генетической предрасположенности не было в 66.51\% (433/655
случаев ) {[}95\% ДИ 0.62;0.7{]} и в 33.49\% было (218/655 случаев )
{[}95\% ДИ 0.3;0.37{]}.

В группе D при исследовании мутации BRCA не было выявлено 2.29\% (15/655
случаев ) {[}95\% ДИ 0.01;0.04{]}, мутация BRCA2 в 0.61\% (4/655 случаев
) {[}95\% ДИ 0;0.02{]}, образование не визуализировалось на ММГ было
1.68\% (11/655 случаев ) {[}95\% ДИ 0.01;0.03{]}, ложноположительный
результат был в 0.61\% (4/655 случаев ) {[}95\% ДИ 0;0.02{]}.

В группе D сторона поражения при осмотре была правая в 12.98\% (85/655
случаев ) {[}95\% ДИ 0.11;0.16{]}, левая в 14.05\% (92/655 случаев )
{[}95\% ДИ 0.12;0.17{]}, поражение обеих сторон было в 6.26\% (41/655
случаев ) {[}95\% ДИ 0.05;0.08{]}, в остальных случаях образование не
пальпировалось.

В группе D кожных симптомов при осмотре не было в 98.63\% (646/655
случаев ) {[}95\% ДИ 0.97;0.99{]}, отек был 0.46\% (3/655 случаев )
{[}95\% ДИ 0;0.01{]}, гиперемия была в 0.92\% (6/655 случаев ) {[}95\%
ДИ 0;0.02{]}.

В группе D симптома втягивания соска при осмотре нет было в 96.49\%
(632/655 случаев ) {[}95\% ДИ 0.95;0.98{]}, симптом был в 3.51\% (23/655
случаев ) {[}95\% ДИ 0.02;0.05{]}. В группе D выделения из соска при
осмотре не был в 98.32\% (644/655 случаев ) {[}95\% ДИ 0.97;0.99{]},
светлые выделения не были 1.22\% (8/655 случаев ) {[}95\% ДИ
0.01;0.02{]}, кровянистые были в 0.46\% (3/655 случаев ) {[}95\% ДИ
0;0.01{]}.

В группе D тип плотности по ACR при осмотре был «С» 88.55\% (580/655
случаев ) {[}95\% ДИ 0.86;0.91{]}, тип «D» былл 11.45\% (75/655 случаев
) {[}95\% ДИ 0.09;0.14{]}.

В группе D1 квадрант локализации верхне-наружный было в13.44\% (88/655
случаев ) {[}95\% ДИ 0.11;0.16{]}, верхне-внутренний в 7.48\% (49/655
случаев ) {[}95\% ДИ 0.06;0.1{]}, нижне-наружный в 2.75\% (18/655
случаев ) {[}95\% ДИ 0.02;0.04{]}, нижне-внутренний в 1.07\% (7/655
случаев ) {[}95\% ДИ 0;0.02{]}, граница наружного в 3.36\% (22/655
случаев ) {[}95\% ДИ 0.02;0.05{]}, граница верхнего в 1.53\% (10/655
случаев ) {[}95\% ДИ 0.01;0.03{]}, граница нижних в 1.07\% (7/655
случаев ) {[}95\% ДИ 0;0.02{]}, вся железа в 1.07\% (7/655 случаев )
{[}95\% ДИ 0;0.02{]}., центральный квадрант было 1.07\% (7/655 случаев )
{[}95\% ДИ 0;0.02{]}.

\hypertarget{ux441ux442ux430ux442ux438ux441ux442ux438ux447ux435ux441ux43aux438ux439-ux430ux43dux430ux43bux438ux437-ux434ux430ux43dux43dux44bux445}{%
\subsection*{2.6 Статистический анализ
данных}\label{ux441ux442ux430ux442ux438ux441ux442ux438ux447ux435ux441ux43aux438ux439-ux430ux43dux430ux43bux438ux437-ux434ux430ux43dux43dux44bux445}}
\addcontentsline{toc}{subsection}{2.6 Статистический анализ данных}

\hypertarget{refs}{}
\begin{CSLReferences}{1}{0}
\leavevmode\vadjust pre{\hypertarget{ref-mandelson2000breast}{}}%
Mandelson, Margaret T, Nina Oestreicher, Peggy L Porter, Donna White,
Charles A Finder, Stephen H Taplin, and Emily White. 2000. {``Breast
Density as a Predictor of Mammographic Detection: Comparison of
Interval-and Screen-Detected Cancers.''} \emph{Journal of the National
Cancer Institute} 92 (13): 1081--87.

\leavevmode\vadjust pre{\hypertarget{ref-narayan2020}{}}%
Narayan, Anand K., Christoph I. Lee, and Constance D. Lehman. 2020.
{``Screening for Breast Cancer.''} \emph{Medical Clinics of North
America} 104 (6): 1007--21.
\url{https://doi.org/10.1016/j.mcna.2020.08.003}.

\leavevmode\vadjust pre{\hypertarget{ref-spak2017}{}}%
Spak, D. A., J. S. Plaxco, L. Santiago, M. J. Dryden, and B. E. Dogan.
2017. {``BI-RADS ® Fifth Edition: A Summary of Changes.''}
\emph{Diagnostic and Interventional Imaging} 98 (3): 179--90.
\url{https://doi.org/10.1016/j.diii.2017.01.001}.

\end{CSLReferences}

\end{document}
